\documentclass[11pt,a4paper]{article}
\usepackage[margin=1in]{geometry}
\usepackage{amsmath,amssymb,booktabs,tabularx,hyperref,microtype}
\hypersetup{colorlinks=true,linkcolor=blue,citecolor=blue,urlcolor=blue,
  pdftitle={Expanded Appendix: CBMR group and covariate inference},
  pdfsubject={Standalone mathematical appendix for Yu et al., Imaging Neuroscience 3 (2025)}}
\emergencystretch=3em
\newcommand{\R}{\mathbb R}
\newcommand{\E}{\operatorname{E}}
\newcommand{\Var}{\operatorname{Var}}
\newcommand{\Cov}{\operatorname{Cov}}
\newcommand{\diag}{\operatorname{diag}}
\newcommand{\blkdiag}{\operatorname{blkdiag}}
\newcommand{\rank}{\operatorname{rank}}
\newcommand{\softplus}{\operatorname{softplus}}
\newcolumntype{T}{>{\raggedright\arraybackslash}X}
\numberwithin{equation}{section}
\setcounter{tocdepth}{2}
\setcounter{secnumdepth}{2}
\title{CBMR: Coordinate-based meta-regression for group and covariate inference\\
  \large Standalone expanded mathematical appendix to Yu et al. (2025)}
\author{}
\date{15 September 2026}

\begin{document}
\maketitle
\begin{abstract}
This editable appendix develops the model and inference calculations for
Yu et al., \emph{Imaging Neuroscience} 3 (2025),
DOI \href{https://doi.org/10.1162/IMAG.a.1057}{10.1162/IMAG.a.1057}.
It concerns group-specific spatial baselines and shared, globally constant
study-covariate effects, exactly the regression setup of that article.
Poisson compression, moment-matched negative-binomial aggregation,
study-clustered negative-binomial integration, complete likelihood
derivatives, structured covariance, and group and covariate tests are
derived step by step. Exact algebra, working-distribution approximations,
and asymptotic inference are distinguished throughout. The paper's
row-vector covariate notation and selected-group/covariate contrasts
are retained, with new helper symbols explicitly distinguished.
Qualifications
address covariate information lost by aggregation, dispersion held fixed
versus jointly estimated, identifiability of a totals-only model with
free group scales, and the uncertainty interpretation of smoothing.
These are self-contained handwritten derivations for manual review, not
full formal-proof certification of the article or of this document.
\end{abstract}
\noindent\textbf{Scope.}
The sole scientific reference is the supplied
\texttt{docs/imag.a.1057.pdf}. This is an explanatory companion, not the
article's original supplementary appendix. Repository derivation notes
provided algebraic provenance; no other manuscript's model is required.
All equations below are numbered independently of the paper.

\tableofcontents
\clearpage

\section{Paper-specific setup and correspondence}\label{sec:setup}

\subsection{Indices, dimensions, and the regression model}
Let $\mathcal I_g$ be the set of studies in group $g$,
$M_g=|\mathcal I_g|$, and $\sum_gM_g=M$.
There are $N$ voxels in a fixed analysis mask and $G$ disjoint groups.
For study $i$, $g(i)$ denotes its group. The observed count at voxel $j$
is $Y_{ij}$, with positive mean $\mu_{ij}$.
The paper notes that these counts are usually zero or one.
Poisson and NB distributions nevertheless assign mass to larger counts:
they are count-model approximations, not normalized Bernoulli models.

\begin{center}
\begin{tabularx}{\textwidth}{l l T}
\toprule
Symbol & Dimension & Meaning\\
\midrule
$X,\ x_j$ & $N\times P,\ P\times1$ & Cubic B-spline design;
row $j$ of $X$ is $x_j^\top$\\
$\beta_g$ & $P\times1$ & Spatial baseline coefficients for group $g$\\
$Z,\ Z_i$ & $M\times R,\ 1\times R$ & Study design and its
row $i$, in the paper's row-vector convention\\
$\gamma$ & $R\times1$ & Shared, globally constant covariate coefficients\\
$\vartheta$ & $(GP+R)\times1$ & Stacked regression coefficients only\\
$\alpha_g$ & scalar & Original group NB overdispersion\\
$r'_{gj},\ p'_{gj},\ \alpha'_g$ & scalars & Paper's matched NB
shape, count probability, and effective dispersion\\
$\nu_g$ & scalar & Clustered gamma shape and rate, $1/\alpha_g$\\
$H,\ I$ & $(GP+R)\times(GP+R)$ & Observed and expected information\\
\bottomrule
\end{tabularx}
\end{center}
The symbol $R$ always counts global covariates; it never denotes a shape.
The paper's single-group and multi-group setups are captured by
\begin{align}
 \eta_{ij}=\log\mu_{ij}
 &=x_j^\top\beta_{g(i)}+{Z_i}\gamma,
 &\vartheta&=(\beta_1^\top,\ldots,\beta_G^\top,\gamma^\top)^\top,
 \label{eq:predictor}\\
 \mu_{gj}^X&=\exp(x_j^\top\beta_g),
 &\mu_i^Z&=\exp({Z_i}\gamma),&
 \mu_{ij}&=\mu_{g(i)j}^X\mu_i^Z. \label{eq:separable}
\end{align}
The superscripts $X$ and $Z$ label the paper's spatial and study
mean factors; they are not powers. In particular $Z_i\gamma$ is
a scalar, $\nabla_\gamma\mu_i^Z=\mu_i^Z Z_i^\top$ is a column,
and $\nabla_\gamma^2\mu_i^Z=\mu_i^Z Z_i^\top Z_i$ is an $R\times R$
matrix. Transposing $Z_i$ in the predictor would reverse the
paper's row convention.
Thus $\partial\eta_{ij}/\partial Z_{ir}=\gamma_r$ is constant over voxels.
A spatially varying baseline is \emph{not} a spatially varying moderator
effect. No coefficient-map extension for moderators is developed here.
The single-group stacked predictor is
$(\mathbf1_M\otimes X)\beta+(Z\otimes\mathbf1_N)\gamma$,
recovering paper Equation~2 with studies stacked before voxels.

\subsection{Native paper symbols and explicitly new helper quantities}
The paper uses $\theta$ for the complete parameter vector in its
likelihoods, and $\theta_j$ for a different object in voxelwise inference.
We keep those meanings. The regression-only vector $\vartheta$
is an additional helper, not a redefinition of $\theta$:
\begin{align}
 \theta_{\rm P}&=\vartheta,\\
 \theta_{\rm NB}
 &=(\beta_1^\top,\ldots,\beta_G^\top,
                 \alpha'_1,\ldots,\alpha'_G,\gamma^\top)^\top,\\
 \theta_{\rm CNB}
 &=(\beta_1^\top,\ldots,\beta_G^\top,
                 \alpha_1,\ldots,\alpha_G,\gamma^\top)^\top .
\end{align}
These match the parameter lists in paper Equations 6, 9, and 11.
An alternative original-dispersion coordinate for the matched NB
replaces $\alpha'_g$ by $\alpha_g$ using Equation~8; its fixed-coordinate
derivatives are distinguished below.

\begin{center}\small
\begin{tabularx}{\textwidth}{lT}
\toprule
Paper symbol & Definition and use here\\
\midrule
$Y_g,\ Y_{gj}$ & Group--voxel counts:
$Y_{gj}=\sum_{i\in\mathcal I_g}Y_{ij}$, with $Y_g\in\R^N$\\
$Y_t,\ Y_{ti}$ & Study totals:
$Y_{ti}=\sum_jY_{ij}$, with $Y_t\in\R^M$\\
$\mu_{gj},\ \mu_{ti}$ & Means of those two totals:
$\mu_{gj}=\sum_{i\in\mathcal I_g}\mu_{ij}$ and
$\mu_{ti}=\sum_j\mu_{ij}$\\
$\mu_g^X,\ \mu^Z$ & Spatial and study mean factors, not the
mean of a group total or the mean of a study total\\
$r'_{gj},\ p'_{gj}$ & Matched NB parameters; derivation below shows
$r'_{gj}=r'_g$ for the paper's separable predictor\\
$\theta_j,\ V_j$ & The $S$ selected groups' voxelwise values and
their $S\times S$ covariance, as in paper Equation~12\\
$C,\ m,\ S$ & Group contrast $C\in\R^{m\times S}$, with $m$
tested constraints and $S$ involved groups\\
$C_\gamma,\ s$ & Covariate contrast $C_\gamma\in\R^{m\times s}$
on $s$ selected covariates, as in paper Equation~14\\
$W_j$ & Signed scalar Wald statistic of paper Equation~13;
its square is the corresponding chi-square statistic\\
\bottomrule
\end{tabularx}
\end{center}
The occurrence $\mu_{g(i)j}$ in paper Equation~5 denotes study $i$'s
cell mean with its group made explicit; we write it as $\mu_{ij}$.
Without the study index, $\mu_{gj}=\E Y_{gj}$ is instead the
\emph{aggregate} mean in Equations~6--9. This distinction is essential.
The paper's $S$ counts selected groups, so we write the Schur complement
as $\mathcal S$, never plain $S$.
The helper $\nu_g=1/\alpha_g$ denotes clustered gamma shape and rate,
while $\lambda_i$ remains the paper's random study multiplier.
New smoothing strengths are denoted $\tau_g$, to keep them distinct
from that random multiplier.
The symbols $H,I,U,\Psi,A,s_1,s_2,h_g,T_g$ and calligraphic
coefficient-contrast matrices below are explicitly defined algebraic
helpers; they are not additional fitted models.
Lowercase $y$ in a scalar mass $p(y)$ is an observed integer argument.

\subsection{Stacking the grouped design one row at a time}
Define $L_{ig}=\mathbf1\{g(i)=g\}$, so $L\in\R^{M\times G}$.
Row $(i,j)$ of $L\otimes X$ is $e_{g(i)}^\top\otimes x_j^\top$.
Row $(i,j)$ of $Z\otimes\mathbf1_N$ is $Z_i$.
Therefore, with study index varying slowest,
\begin{equation}
 \mathcal X=[\,L\otimes X\quad Z\otimes\mathbf1_N\,],
 \qquad \eta=\mathcal X\vartheta,\qquad
 \dim(\mathcal X)=MN\times(GP+R).
\end{equation}
The $X$ in this formula is still the paper's $N\times P$ spatial
matrix; $\mathcal X$ is a new name for the much larger combined
design, which need not be constructed numerically.
For $G=1$, $L=\mathbf1_M$, giving the single-group design of
paper Equation~2. This verifies the coefficient ordering used in
all later block matrices.

\subsection{Group-specific constant effects allowed by the paper}
Paper Section 1.2.2 describes both independent group fits and a
single fit with group-supported covariate columns.
For a covariate row $Z_i\in\R^{1\times R_0}$ whose effects are all
group-specific, construct
\[
 Z^{\rm sep}_{i,(g,r)}
   =\mathbf1\{g(i)=g\}Z_{ir},\qquad
 \gamma^{\rm sep}=(\gamma_1^\top,\ldots,\gamma_G^\top)^\top.
\]
Then $Z_i^{\rm sep}\gamma^{\rm sep}=Z_i\gamma_{g(i)}$, and
\[
 \eta_{ij}=x_j^\top\beta_{g(i)}+Z_i\gamma_{g(i)}.
\]
For two groups, two studies per group, and a single covariate with
values $a_1,\ldots,a_4$, the new study design is
\[
 Z^{\rm sep}=
 \begin{pmatrix}a_1&0\\a_2&0\\0&a_3\\0&a_4\end{pmatrix}.
\]
There are still scalar moderator effects constant over voxels.
This construction does not introduce spatial coefficient maps.
With only group-specific parameters, independent group data, and
separate group penalties, the likelihood factors in parameter
blocks $\zeta_g=(\beta_g^\top,\gamma_g^\top)^\top$.
Cross derivatives in $\zeta_g,\zeta_h$ vanish for $g\ne h$, and
the full information and covariance are block diagonal after
ordering by $\zeta_g$. This is why the two fitting strategies agree
under those conditions.

For a mixture of common and group-specific covariates, write
\[
 \eta_{ij}=x_j^\top\beta_{g(i)}
            +Z_i^{\rm common}\gamma^{\rm common}
            +Z_i^{\rm specific}\gamma_{g(i)}^{\rm specific}.
\]
Combine $(\beta_g,\gamma_g^{\rm specific})$ into each group block.
The common coefficients form the remaining border, so the later
Schur formulas apply with these enlarged group blocks.
Groups are not independent after estimating that shared border.
Constant-within-group covariates can still be confounded with
spatial intercepts; encoding them in separate columns does not
by itself create identifiability.

\subsection{Correspondence with the published equations}
\begin{center}
\begin{tabularx}{\textwidth}{l T T}
\toprule
Paper location & Subject & Expansion here\\
\midrule
Equations 1--5 & Single/multiple groups, shared covariates,
separable mean & Sections \ref{sec:setup}--\ref{sec:chain}\\
Equation 6 & Poisson likelihood and compression &
Section \ref{sec:poisson}, including the allocation qualification\\
Equations 7--9 & NB moment matching and aggregate likelihood &
Sections \ref{sec:cellnb}--\ref{sec:aggregate}\\
Equations 10--11 & Within-study dependence and clustered likelihood &
Section \ref{sec:cluster}\\
Algorithm 1 & Alternating dispersion and regression updates &
Section \ref{sec:optimization}\\
Equations 12--14 & Group-map and global-covariate Wald inference &
Sections \ref{sec:covariance}--\ref{sec:inference}\\
Section 2.1.4 & Bootstrap and numerical caveats &
Sections \ref{sec:bootstrap}--\ref{sec:numerics}\\
\bottomrule
\end{tabularx}
\end{center}
Paper Section 2.1.2 explicitly excludes quasi-Poisson from its fitted
pipeline, despite introductory discussion of earlier work. Accordingly
there is no fourth fitted quasi-Poisson family in this appendix.
The independent-cell NB below is the generative starting point for
the paper's moment-matched NB, not a claim that an additional
independent-cell fitting option was evaluated in that article.

\subsection{Identification and the meaning of information}
Covariates are preprocessed and centered as described in the paper.
Centering improves interpretation, but does not by itself prove rank.
If $Xc=\mathbf1_N$ and ${Z_i} d=b_{g(i)}$ for every study, then
\begin{equation}
 \beta_g\mapsto\beta_g-tb_gc,\qquad
 \gamma\mapsto\gamma+td
 \quad\Longrightarrow\quad \eta_{ij}\text{ unchanged}.
 \label{eq:confounding}
\end{equation}
An intercept in $Z$, a group-constant covariate, redundant spline columns,
or combinations of such directions require constraints or removal.
Constraints must also be respected when forming contrasts and covariance.
Full rank, finite fitted means, and an interior finite optimum are separate
requirements; an all-zero region can drive coefficients towards a boundary.

We write $\ell$ for log likelihood, $L=-\ell$ for negative log likelihood,
$U=\nabla\ell$ for score, and
\begin{equation}
 H=-\nabla^2\ell=\nabla^2L,\qquad I=\E_\vartheta H.
 \label{eq:information}
\end{equation}
Dispersion is held fixed in regression derivatives unless stated otherwise.
Under a correctly specified regular likelihood, differentiation under
the expectation gives $\E U=0$ and $I=\Var(U)$.
Those are distributional facts with regularity conditions.
Inverting a matrix is exact algebra; identifying its inverse with estimator
covariance is normally a first-order asymptotic approximation.
A working likelihood may have exact derivatives without being the exact
distribution of the data, and its inverse Hessian need not give robust
sampling uncertainty.

\section{Affine differentiation and group information blocks}\label{sec:chain}

\subsection{From a scalar mean to a coefficient Hessian}
Let $e_g$ be column $g$ of the $G$-dimensional identity and define
\begin{equation}
 b_{ij}=\begin{pmatrix}e_{g(i)}\otimes x_j\\{Z_i^\top}\end{pmatrix},
 \qquad \eta_{ij}=b_{ij}^\top\vartheta .
\end{equation}
The cell-design vector $b_{ij}$ is not the spatial row $x_j$.
For a single cell, suppress its indices and let $\mu=e^\eta$.
Direct differentiation gives
\begin{align}
 \partial_a\eta&=b_a,& \partial_{ab}\eta&=0,&
 \partial_a\mu&=\mu b_a,& \partial_{ab}\mu&=\mu b_ab_b.
\end{align}
For any twice differentiable scalar function $\ell(\mu)$,
\begin{align}
 \partial_a\ell&=\ell_\mu\mu b_a=\ell_\eta b_a,\\
 \partial_{ab}\ell
 &=\ell_{\mu\mu}\mu^2b_ab_b+\ell_\mu\mu b_ab_b
 =\ell_{\eta\eta}b_ab_b. \label{eq:affine-chain}
\end{align}
The $\ell_\mu\mu$ term is essential: differentiating only the outer
function with respect to $\mu$ is not the log-link Hessian.
Equivalently, $d\ell=\ell_\eta b^\top d\vartheta$ and
$d^2\ell=\ell_{\eta\eta}(b^\top d\vartheta)^2$.

\subsection{Independent factors and spatial support}
When cells are independent, write
$s_{ij}=\partial\ell_{ij}/\partial\eta_{ij}$ and
$w_{ij}=-\partial^2\ell_{ij}/\partial\eta_{ij}^2$.
Summing \eqref{eq:affine-chain} yields
\begin{equation}
 U=\sum_{i,j}b_{ij}s_{ij},\qquad
 H=\sum_{i,j}w_{ij}b_{ij}b_{ij}^\top. \label{eq:cell-assembly}
\end{equation}
Its individual blocks are
\begin{align}
 H_{\beta_g\beta_g}
 &=\sum_{i\in\mathcal I_g}\sum_jw_{ij}x_jx_j^\top,\\
 H_{\beta_g\beta_h}&=0\quad(g\ne h),\\
 H_{\beta_g\gamma}
 &=\sum_{i\in\mathcal I_g}\sum_jw_{ij}x_j{Z_i},\\
 H_{\gamma\gamma}&=\sum_{i,j}w_{ij}{Z_i^\top}{Z_i}. \label{eq:group-blocks}
\end{align}
The zero cross-group block follows from mutually exclusive membership,
not from treating the shared $\gamma$ as known. Its inverse can have
nonzero cross-group blocks.

For a study likelihood coupling voxels, let $B_i$ have rows $b_{ij}^\top$,
$s_i=\nabla_{\eta_i}\ell_i$, and $W_i=-\nabla_{\eta_i}^2\ell_i$.
Then $U_i=B_i^\top s_i$ and $H_i=B_i^\top W_iB_i$.
The matrix $W_i$ need not be diagonal.
Yet a study involves only one $\beta_g$, so the same cross-group
structural zeros remain. For aggregate NB, the factor is instead a
group--voxel total; its full composite derivatives must be used.

\section{Poisson likelihood and exact two-marginal compression}\label{sec:poisson}

\subsection{Probability mass, score, and curvature}
For independent cells,
\begin{align}
 \Pr(Y=y\mid\mu)&=e^{-\mu}\mu^y/y!,\\
 \ell_{ij}&=Y_{ij}\eta_{ij}-e^{\eta_{ij}}-\log(Y_{ij}!),\\
 s_{ij}&=Y_{ij}-\mu_{ij},&
 w_{ij}&=\mu_{ij}. \label{eq:poisson-cell}
\end{align}
The score is the residual because $\partial_\eta e^\eta=e^\eta$;
its derivative is $-\mu$ because the observed count is fixed.
Consequently
\begin{equation}
 U=\sum_{i,j}b_{ij}(Y_{ij}-\mu_{ij}),\qquad
 H=I=\sum_{i,j}\mu_{ij}b_{ij}b_{ij}^\top.
\end{equation}
For any $a$, $a^\top Ha=\sum_{ij}\mu_{ij}(b_{ij}^\top a)^2\geq0$.
At finite positive means it is strictly positive for nonzero $a$
exactly when the stacked regression design has no null direction.
This proves convexity of the negative log likelihood, not existence
of a finite minimizer in a data-poor design.

\subsection{Compress the full likelihood without discarding covariate data}
Define the two observed margins and two mean totals
\begin{equation}
 Y_{gj}=\sum_{i\in\mathcal I_g}Y_{ij},\quad
 Y_{ti}=\sum_jY_{ij},\quad
 h_g=\sum_j\mu_{gj}^X,\quad T_g=\sum_{i\in\mathcal I_g}\mu_i^Z.
 \label{eq:margins}
\end{equation}
Separate each part of the cell sum:
\begin{align}
 \sum_{i,j}Y_{ij}x_j^\top\beta_{g(i)}
 &=\sum_g Y_g^\top X\beta_g,\\
 \sum_{i,j}Y_{ij}{Z_i}\gamma&=\sum_i Y_{ti} {Z_i}\gamma,\\
 \sum_{i,j}\mu_{ij}
 &=\sum_g\left(\sum_j\mu_{gj}^X\right)
            \left(\sum_{i\in\mathcal I_g}\mu_i^Z\right)=\sum_gh_gT_g .
\end{align}
Thus, up to a parameter-independent data term,
\begin{equation}
 \ell_{\rm full}
 =\sum_g Y_g^\top X\beta_g+\sum_i Y_{ti} {Z_i}\gamma
       -\sum_gh_gT_g+c(Y). \label{eq:poisson-compressed}
\end{equation}
Both spatial totals $Y_g$ and study totals $Y_{ti}$ are retained.
This is an exact compression of the \emph{full} Poisson likelihood.
The factorial data term may be dropped for parameter optimization,
but likelihood comparisons requiring absolute normalizations must
keep track of which observation space is being used.

\subsection{Differentiate the compressed form explicitly}
Define
\begin{align}
 a_g&=X^\top \mu_g^X,&
 B_g&=X^\top\diag(\mu_g^X)X,\\
 t_g^{(z)}&=\sum_{i\in\mathcal I_g}\mu_i^Z{Z_i^\top},&
 V_g&=\sum_{i\in\mathcal I_g}\mu_i^Z{Z_i^\top}{Z_i} .
\end{align}
Then $\nabla_{\beta_g}h_g=a_g$, $\nabla_{\beta_g}^2h_g=B_g$,
$\nabla_\gamma T_g=t_g^{(z)}$, and $\nabla_\gamma^2T_g=V_g$.
Differentiating \eqref{eq:poisson-compressed} once gives
\begin{equation}
 U_{\beta_g}=X^\top Y_g-T_ga_g,\qquad
 U_\gamma=\sum_i Y_{ti}{Z_i^\top}-\sum_gh_gt_g^{(z)} .
\end{equation}
Differentiating again and negating gives
\begin{equation}
 D_g=T_gB_g,\qquad F_g=a_g(t_g^{(z)})^\top,\qquad
 E=\sum_gh_gV_g . \label{eq:poisson-blocks}
\end{equation}
Here $D_g=H_{\beta_g\beta_g}$, $F_g=H_{\beta_g\gamma}$,
and $E=H_{\gamma\gamma}$; $E$ is a matrix, whereas $\E$ is expectation.
These formulas reproduce the cell sums without an $M\times N$ mean array.
For example the $P\times R$ cross block comes from differentiating
$-T_ga_g$ with respect to $\gamma$, not from an independence assumption.

\subsection{The two expressions in paper Equation 6 need a qualification}
Poisson closure says $Y_{gj}\sim\operatorname{Poisson}(\mu_{gj}^XT_g)$.
The likelihood using \emph{only} these totals is
\begin{equation}
 \ell_{\rm totals}
 =\sum_g Y_g^\top X\beta_g+\sum_g Y_{g\cdot}\log T_g-\sum_gh_gT_g+c_{\rm tot}(Y),
 \qquad Y_{g\cdot}=\sum_{i\in\mathcal I_g}Y_{ti} .
\end{equation}
This corresponds to the totals-only first expression in paper Equation~6.
The paper's later expression retaining study totals corresponds to
\eqref{eq:poisson-compressed}. Subtraction exposes the distinction:
\begin{equation}
 \ell_{\rm full}-\ell_{\rm totals}
 =\sum_g\sum_{i\in\mathcal I_g}Y_{ti}\log(\mu_i^Z/T_g)
       +c(Y)-c_{\rm tot}(Y). \label{eq:allocation}
\end{equation}
Conditional on a voxel total, Poisson allocations across studies are
multinomial with probabilities $\pi_{ig}=\mu_i^Z/T_g$.
Indeed, dividing the product cell mass by the total mass gives
$Y_{gj}!\prod_i\pi_{ig}^{Y_{ij}}/\prod_iY_{ij}!$.
Summing its log over voxels gives exactly \eqref{eq:allocation}.

For fixed $\gamma$, this difference is constant in $\beta$.
Without estimated covariates it is also constant in the remaining
regression parameters. In general it is not constant in estimated $\gamma$.
Writing $\bar z_g=t_g^{(z)}/T_g$,
\begin{align}
 U_{{\rm alloc},g}&=\sum_{i\in\mathcal I_g}Y_{ti}{Z_i^\top}-Y_{g\cdot}\bar z_g,\\
 H_{{\rm alloc},g}
 &=Y_{g\cdot}\{V_g/T_g-\bar z_g\bar z_g^\top\}
 =Y_{g\cdot}\Cov_{\pi_g}({Z_i^\top})\succeq0 .
\end{align}
Therefore Poisson closure does not justify blanket equivalence of
totals-only and full likelihoods when covariate effects are estimated.
The Poisson limit of the aggregate NB likelihood below is the totals-only
likelihood, not automatically the full two-margin expression.

\section{Independent NB cells as the basis of moment matching}\label{sec:cellnb}

\subsection{NB2 convention and the underlying probability mass}
Within group $g$, put $r=1/\alpha_g>0$.
An independent gamma multiplier with shape and rate $r$ has mean one
and variance $\alpha_g$. Mixing a Poisson mean $\mu$ by that multiplier
gives
\begin{equation}
 p(y\mid\mu,r)=\frac{\Gamma(y+r)}{\Gamma(r)\Gamma(y+1)}
       \left(\frac{r}{r+\mu}\right)^r
       \left(\frac{\mu}{r+\mu}\right)^y,
 \qquad \Var(Y)=\mu+\alpha_g\mu^2.
 \label{eq:nb-pmf}
\end{equation}
The cells in this construction have independent multipliers, unlike
the common study multiplier in Section~\ref{sec:cluster}.
At fixed $r$, the scalar log likelihood is
\begin{equation}
 \ell=\log\Gamma(y+r)-\log\Gamma(r)-\log\Gamma(y+1)
       +r\log r+y\eta-(y+r)\log(r+e^\eta).
\end{equation}
Retain the gamma terms if dispersion is differentiated.
They are constant only during the regression step at fixed dispersion.

\subsection{Scalar derivatives, information, and their limitations}
First combine the score numerator:
\begin{equation}
 \ell_\eta=y-(y+r)\frac{\mu}{r+\mu}
 =\frac{r(y-\mu)}{r+\mu}=\frac{y-\mu}{1+\alpha_g\mu}.
\end{equation}
The quotient rule now gives
\begin{align}
 \ell_{\eta\eta}
 &=\frac{-\mu(1+\alpha_g\mu)-\alpha_g\mu(y-\mu)}
                  {(1+\alpha_g\mu)^2}
 =-\frac{\mu(1+\alpha_gy)}{(1+\alpha_g\mu)^2},\\
 w^{\rm obs}&=\frac{\mu(1+\alpha_gy)}{(1+\alpha_g\mu)^2},&
 w^{\rm F}&=\E w^{\rm obs}=\frac{\mu}{1+\alpha_g\mu}. \label{eq:nb-weights}
\end{align}
Also $\Var(\ell_\eta)=(\mu+\alpha_g\mu^2)/(1+\alpha_g\mu)^2=w^{\rm F}$.
Substitution in \eqref{eq:cell-assembly}--\eqref{eq:group-blocks}
gives all underlying independent-cell regression blocks.
Their observed curvature is nonnegative at fixed $\alpha_g$,
and both weights tend to $\mu$ as $\alpha_g\downarrow0$.
The denominators contain cell means; simply replacing $Y_{ij}$ by
$Y_{gj}$ does not produce the paper's moment-matched likelihood.
This cell-level derivation identifies the original dispersion convention
and clarifies why Poisson compression cannot be reused unaltered.

\subsection{Shape derivatives for an optional joint cell-model calculation}
Define $\psi=d\log\Gamma/dr$ and $\psi_1=d\psi/dr$.
At fixed $\mu$ the exact derivatives are
\begin{align}
 \ell_r&=\psi(y+r)-\psi(r)+\log r+1-\log(r+\mu)-\frac{y+r}{r+\mu},\\
 \ell_{rr}&=\psi_1(y+r)-\psi_1(r)+\frac1r-\frac1{r+\mu}
                                  -\frac{\mu-y}{(r+\mu)^2},\\
 \ell_{\eta r}&=\frac{\mu(y-\mu)}{(r+\mu)^2}.
\end{align}
Thus $H_{\vartheta r}=-\sum b_{ij}\ell_{\eta r}$ within the relevant group.
Under this correctly specified cell model, $\E\ell_{\eta r}=0$.
Observed cross information need not vanish. Neither that orthogonality
nor fixed-dispersion convexity transfers automatically to the composed
aggregate likelihood with $\gamma$-dependent shape.

\section{Moment-matched aggregate NB: full composite differentiation}
\label{sec:aggregate}

\subsection{Match the mean and variance, not the entire distribution}
Temporarily omit the fixed group label $g$:
$Y_j=Y_{gj}$, $\mu_j=\mu_{gj}$, $\mu_j^X=\mu_{gj}^X$,
$\alpha=\alpha_g$, $\alpha'=\alpha'_g$, and $p'_j=p'_{gj}$.
The study index $i$ still ranges over $\mathcal I_g$.
After deriving voxel-constant shape we write $r'=r'_g=r'_{gj}$.
Let
\begin{equation}
 s_1=\sum_i \mu_i^Z,\qquad s_2=\sum_i(\mu_i^Z)^2,\qquad s_3=\sum_i(\mu_i^Z)^3.
\end{equation}
Independence of the original NB cells gives
\begin{equation}
 {\mu_j}=\E Y_j=\mu_j^Xs_1,\qquad
 \Var(Y_j)=\mu_j^Xs_1+\alpha (\mu_j^X)^2s_2 .
\end{equation}
An NB variable of mean ${\mu_j}$ and shape ${r'}$ has variance
${\mu_j}+{\mu_j}^2/{r'}$. Equating extra variances and cancelling $(\mu_j^X)^2$
gives
\begin{align}
 {r'}&=\frac{s_1^2}{\alpha s_2},&
 A&=\frac{s_1}{\alpha s_2}=\frac{{r'}}{s_1},\\
 {p'}_j&=\frac{{\mu_j}}{{\mu_j}+{r'}}=\frac{\mu_j^X}{\mu_j^X+A},&
 \alpha'&=\frac1{r'}=\frac{\alpha s_2}{s_1^2}.
 \label{eq:matched-parameters}
\end{align}
Our ${p'}_j$ is the probability appearing as ${p'}_j^{Y_j}$ in the NB mass.
These are paper Equations~7--8 specialized to its separable mean.
The effective shape is constant over voxels within a group,
but generally depends on $\gamma$.

\paragraph{Recover the paper's primed parameters explicitly.}
Define the extra variance
$\Delta_{gj}=\alpha_g\sum_{i\in\mathcal I_g}\mu_{ij}^2$.
Here $s_{kg}=\sum_{i\in\mathcal I_g}(\mu_i^Z)^k$ restores the
group index suppressed in $s_k$.
For an NB with count probability $p'$ and shape $r'$,
$\mu=r'p'/(1-p')$ and
$\Var(Y)=r'p'/(1-p')^2$. Subtracting its mean gives
$\Var(Y)-\mu=\mu^2/r'$.
Equating this to $\Delta_{gj}$ and then solving for $p'$ yields
\begin{align}
 r'_{gj}
 &=\frac{\mu_{gj}^2}{\Delta_{gj}}
   =\frac{(\mu_{gj}^X s_{1g})^2}
           {\alpha_g(\mu_{gj}^X)^2s_{2g}}
   =\frac{s_{1g}^2}{\alpha_gs_{2g}}=r'_g,\\
 p'_{gj}
 &=\frac{\mu_{gj}}{\mu_{gj}+r'_{gj}}
   =\frac{\Delta_{gj}}{\mu_{gj}+\Delta_{gj}}
   =\frac{\sum_{i\in\mathcal I_g}\mu_{ij}^2}
           {\alpha_g^{-1}\mu_{gj}+\sum_{i\in\mathcal I_g}\mu_{ij}^2},\\
 \alpha'_g
 &=\frac{\Delta_{gj}}{\mu_{gj}^2}
   =\frac{\alpha_gs_{2g}}{s_{1g}^2}=\frac1{r'_g}.
 \label{eq:paper-prime-derivation}
\end{align}
This is a dimensionally complete derivation of the moment-matching
relations in paper Equations~7--8, including the probability convention.
Only the probability varies over voxels in this separable setup.
The parameter $A_g=r'_g/s_{1g}$ is an additional helper introduced
for differentiation, not a replacement for the paper's $r'_{gj}$
or $\alpha'_g$.

The generating function of an original cell is
$(1+\alpha \mu_j^X\mu_i^Z(1-t))^{-1/\alpha}$.
Its product over studies is generally not one NB generating function.
If all $\mu_i^Z={\mu^Z}$, the probabilities are common and the product is
$(1+\alpha \mu_j^X\mu^Z(1-t))^{-M_g/\alpha}$: the sum is exactly NB.
More generally NB closure requires a common probability parameter.
Here, with common original $\alpha$, equal positive $\mu_i^Z$ supplies it.

To see what matching misses, the NB cumulant generating function is
$K(t)=-\alpha^{-1}\log\{1-\alpha\mu(e^t-1)\}$.
To show all three differentiations, put
$D(t)=1-\alpha\mu(e^t-1)$, so $D'(t)=-\alpha\mu e^t$.
The successive derivatives are
\begin{align}
 K'(t)&=\mu e^t/D(t),\\
 K''(t)&=\mu e^t/D(t)+\alpha\mu^2e^{2t}/D(t)^2,\\
 K'''(t)&=\mu e^t/D(t)+3\alpha\mu^2e^{2t}/D(t)^2
                      +2\alpha^2\mu^3e^{3t}/D(t)^3.
\end{align}
Using $D(0)=1$ gives
$K'''(0)=\mu+3\alpha\mu^2+2\alpha^2\mu^3$.
Add the true independent cumulants and subtract the matched one:
\begin{equation}
 \kappa_3^{\rm true}-\kappa_3^{\rm matched}
 =2\alpha^2(\mu_j^X)^3\left(s_3-\frac{s_2^2}{s_1}\right)\geq0.
 \label{eq:third-cumulant}
\end{equation}
The inequality is Cauchy--Schwarz,
$(\sum_i(\mu_i^Z)^2)^2\leq(\sum_i\mu_i^Z)(\sum_i(\mu_i^Z)^3)$.
Equality for positive $\mu_i^Z$ requires equal $\mu_i^Z$.
Moment matching is a working distribution, not a general convolution identity.

\subsection{Rewrite paper Equation 9 before differentiating}
Set $\eta_j=x_j^\top\beta$, $\mu_j^X=e^{\eta_j}$, $d_j=\mu_j^X+A$.
Ignoring only count-factorial terms, the negative log likelihood is
\begin{align}
 \Psi(\eta,{r'},A)
 ={}&-\sum_j\log\Gamma(Y_j+{r'})+N\log\Gamma({r'})-N{r'}\log A
 \notag\\
 &+\sum_j({r'}+Y_j)\log(\mu_j^X+A)-\sum_jY_j\eta_j .
 \label{eq:Psi}
\end{align}
Indeed, $\log {p'}_j=\eta_j-\log d_j$ and
$\log(1-{p'}_j)=\log A-\log d_j$.
First treat $(\eta,{r'},A)$ as independent coordinates.
Only afterwards substitute ${r'}(\gamma,\alpha)$ and $A(\gamma,\alpha)$.
This order keeps the nonlinear inner-map curvature visible.

\subsection{All outer first and second partial derivatives}
At fixed ${r'},A$,
\begin{align}
 \Psi_{\eta_j}&=({r'}+Y_j)\frac{\mu_j^X}{d_j}-Y_j,\\
 \partial_{\eta_j}(\mu_j^X/d_j)&=\frac{\mu_j^Xd_j-(\mu_j^X)^2}{d_j^2}
                                  =\frac{\mu_j^XA}{d_j^2},\\
 \Psi_{\eta_j\eta_j}&=\frac{({r'}+Y_j)\mu_j^XA}{d_j^2}=w_j,&
 \Psi_{\eta_j\eta_k}&=0\quad(j\ne k),\\
 \Psi_{\eta_j{r'}}&=\frac{\mu_j^X}{d_j},&
 \Psi_{\eta_j A}&=-\frac{({r'}+Y_j)\mu_j^X}{d_j^2}. \label{eq:Psi-spatial}
\end{align}
The shape and scale derivatives are
\begin{align}
 \Psi_{r'}&=-\sum_j\psi(Y_j+{r'})+N\psi({r'})
                         -N\log A+\sum_j\log d_j,\\
 \Psi_A&=-N{r'}/A+\sum_j({r'}+Y_j)/d_j,\\
 \Psi_{{r'}{r'}}&=-\sum_j\psi_1(Y_j+{r'})+N\psi_1({r'}),\\
 \Psi_{{r'} A}&=-N/A+\sum_jd_j^{-1},\\
 \Psi_{AA}&=N{r'}/A^2-\sum_j({r'}+Y_j)/d_j^2. \label{eq:Psi-outer}
\end{align}
For instance $\partial_A(-N{r'}/A)=N{r'}/A^2$,
whereas differentiating each reciprocal $d_j^{-1}$ contributes a minus.
Together \eqref{eq:Psi-spatial}--\eqref{eq:Psi-outer} specify every
first and second outer derivative; omitted off-diagonal $\eta$ entries
are zero, and reversed mixed entries follow by symmetry.

\subsection{Inner gradients and Hessians: the factors two and four}
All derivatives in this subsection are with respect to $\gamma$,
holding the \emph{original} $\alpha$ fixed.
Since $\mu_i^Z=\exp({Z_i}\gamma)$,
\begin{equation}
 \nabla \mu_i^Z=\mu_i^Z{Z_i^\top},\quad \nabla^2\mu_i^Z=\mu_i^Z{Z_i^\top}{Z_i},\quad
 \nabla((\mu_i^Z)^2)=2(\mu_i^Z)^2{Z_i^\top},\quad \nabla^2((\mu_i^Z)^2)=4(\mu_i^Z)^2{Z_i^\top}{Z_i}.
\end{equation}
Define $b_k=\sum_i(\mu_i^Z)^k{Z_i^\top}$ and $C_k=\sum_i(\mu_i^Z)^k{Z_i^\top}{Z_i}$,
for $k=1,2$. Then
\begin{align}
 \nabla s_1&=b_1,&\nabla^2s_1&=C_1,&
 \nabla s_2&=2b_2,&\nabla^2s_2&=4C_2 .
\end{align}
Apply $\nabla^2\log s=s^{-1}\nabla^2s-s^{-2}\nabla s\nabla s^\top$:
\begin{align}
 l&=\nabla\log s_1=b_1/s_1,&
 L_1&=\nabla^2\log s_1=C_1/s_1-b_1b_1^\top/s_1^2,\\
 h&=\nabla\log s_2=2b_2/s_2,&
 L_2&=\nabla^2\log s_2=4C_2/s_2-4b_2b_2^\top/s_2^2 .
\end{align}
With $q=\log{r'}$ and $a=\log A$,
\begin{align}
 q&=2\log s_1-\log s_2-\log\alpha,&
 d_q=\nabla q&=2l-h,&Q_q=\nabla^2q&=2L_1-L_2,\\
 a&=\log s_1-\log s_2-\log\alpha,&
 d_a=\nabla a&=l-h,&Q_a=\nabla^2a&=L_1-L_2 .
 \label{eq:log-inner}
\end{align}
Finally differentiation of $e^q$ and $e^a$ gives
\begin{align}
 r_\gamma:=\nabla{r'}&={r'} d_q,&
 R_\gamma:=\nabla^2{r'}&={r'}(d_qd_q^\top+Q_q),\\
 a_\gamma:=\nabla A&=A d_a,&
 A_\gamma:=\nabla^2A&=A(d_ad_a^\top+Q_a).
 \label{eq:inner}
\end{align}
The capital $R_\gamma$ here is a Hessian, not the scalar number $R$.
The outer-product terms arise from differentiating the exponential
prefactor and must not be discarded.

\subsection{The complete regression score and group blocks}
The spatial score and information are immediately
\begin{equation}
 U_\beta=-X^\top\Psi_\eta,\qquad
 D=X^\top\diag(w)X.
\end{equation}
For a global component $a$, write
$r'_a=\partial r'/\partial\gamma_a$,
$A_a=\partial A/\partial\gamma_a$, and use two component indices
for the corresponding second derivatives. Then
$\partial_{\gamma_a}\Psi=\Psi_{r'}{r'}_a+\Psi_AA_a$.
Differentiating each product with respect to $\gamma_b$ gives
\begin{align}
 \partial_{\gamma_b}(\Psi_{r'}{r'}_a)
 &=\Psi_{{r'}{r'}}{r'}_b{r'}_a+\Psi_{{r'} A}A_b{r'}_a
                                  +\Psi_{r'}{r'}_{ab},\\
 \partial_{\gamma_b}(\Psi_AA_a)
 &=\Psi_{{r'} A}{r'}_bA_a+\Psi_{AA}A_bA_a+\Psi_AA_{ab}.
\end{align}
Thus the group contribution to the global score and Hessian is
\begin{align}
 U_{\gamma,g}&=-\Psi_{r'} r_\gamma-\Psi_Aa_\gamma,\\
 E_g&=\Psi_{{r'}{r'}}r_\gamma r_\gamma^\top
   +\Psi_{{r'} A}(r_\gamma a_\gamma^\top+a_\gamma r_\gamma^\top)
   +\Psi_{AA}a_\gamma a_\gamma^\top
   +\Psi_{r'} R_\gamma+\Psi_A A_\gamma . \label{eq:aggregate-global}
\end{align}
The mixed spatial/global block is
\begin{equation}
 F_g=\sum_jx_j\left\{\frac{\mu_j^X}{d_j}r_\gamma^\top
                 -\frac{({r'}+Y_j)\mu_j^X}{d_j^2}a_\gamma^\top\right\}.
 \label{eq:aggregate-cross}
\end{equation}
Restore the group subscripts on all inner quantities, set
$D=\blkdiag(D_1,\ldots,D_G)$, stack $F_g$, and sum $E=\sum_gE_g$.
Likewise $U_\gamma=\sum_gU_{\gamma,g}$.
This accounts for every regression derivative at fixed original
$\alpha_1,\ldots,\alpha_G$.

In particular $\Psi_{r'} R_\gamma+\Psi_AA_\gamma$ does not vanish
merely because the composite score is zero.
An optimum imposes the available parameter score equations, not
independent optimization over unconstrained ${r'}$ and $A$ for every group.
Discarding these terms is a curvature approximation, not this Hessian.

\subsection{Holding effective dispersion fixed is a different operation}
Paper Equations~8--9 introduce effective dispersion $\alpha'$.
Its use as an optimization coordinate must be distinguished from
holding original $\alpha$ fixed in \eqref{eq:log-inner}.
If $\alpha'$ is held fixed, then
\begin{equation}
 q=-\log\alpha',\quad a=q-\log s_1,\quad
 d_q=0,\ Q_q=0,\ d_a=-l,\ Q_a=-L_1 .
 \label{eq:effective-fixed}
\end{equation}
These derivatives differ from \eqref{eq:log-inner}.
The matched mean is still $\mu_j^Xs_1$, but the shape no longer changes with
$\gamma$ during that partial update.
Because $\alpha'=\alpha s_2/s_1^2$, a joint, unconstrained positive
dispersion parameterization can describe the same matched distributions.
However, a fixed-dispersion regression slice, an alternating update path,
and a conditional information block depend on which coordinate is held.
Joint covariance of the common regression coordinates is invariant under
a smooth nuisance reparameterization when transformed consistently
on identified coordinates where that covariance exists.

\paragraph{Direct Hessian in the paper's effective-dispersion coordinate.}
At fixed $\alpha'$, the aggregate log mean is
$\widetilde\eta_j=x_j^\top\beta+\log s_1$, not an affine function
of $\gamma$. Define
\[
 \widetilde s_j=\frac{Y_j-\mu_j}{1+\alpha'\mu_j},\qquad
 w_j=\frac{\mu_j(1+\alpha'Y_j)}{(1+\alpha'\mu_j)^2}.
\]
The tilde distinguishes these scores from the study moment sums $s_1,s_2$;
write $\widetilde s=(\widetilde s_j)_j$ and $w=(w_j)_j$ below.
The derivatives of the predictor are
$\nabla_\beta\widetilde\eta_j=x_j$,
$\nabla_\gamma\widetilde\eta_j=l$, and
$\nabla_\gamma^2\widetilde\eta_j=L_1$.
Applying the full composite rule gives the group contributions
\begin{align}
 U_\beta&=X^\top\widetilde s,& U_{\gamma,g}&=l\sum_j\widetilde s_j,\\
 D_g&=X^\top\diag(w)X,& F_g&=(X^\top w)l^\top,\\
 E_g&=\left(\sum_jw_j\right)ll^\top-\left(\sum_j\widetilde s_j\right)L_1.
 \label{eq:effective-direct}
\end{align}
In the last line, $-\sum_j\widetilde s_jL_1$ comes from
$-\ell_{\widetilde\eta_j}\nabla_\gamma^2\widetilde\eta_j$.
Dropping it away from a zero spatial-intercept score is not the
observed Hessian. These formulas are the direct reduction of the
outer/inner calculation with \eqref{eq:effective-fixed}.

\subsection{An identification consequence of using only matched totals}
\label{sec:aggregate-identification}
Suppose the spatial basis contains a freely fitted intercept
$Xc_B=\mathbf1_N$. For any two covariate vectors $\gamma$ and
$\widetilde\gamma$, set
\[
 \widetilde\beta_g=\beta_g+
   \log\!\left\{\frac{s_{1g}(\gamma)}
                      {s_{1g}(\widetilde\gamma)}\right\}c_B,
 \qquad \widetilde\alpha'_g=\alpha'_g.
\]
Then
\begin{align}
 \widetilde\mu_{gj}
 &=\exp(x_j^\top\widetilde\beta_g)s_{1g}(\widetilde\gamma)
   =\exp(x_j^\top\beta_g)s_{1g}(\gamma)=\mu_{gj},\\
 \widetilde r'_{gj}&=1/\widetilde\alpha'_g=r'_{gj}.
\end{align}
Every factor in the matched-totals likelihood is unchanged.
Consequently, these totals alone do not identify the shared
$\gamma$ separately from freely fitted group intercepts, even
when the study design has full rank and varies within groups.
This invariance holds with effective dispersion fixed or freely
estimated. With original dispersion freely estimated, the same
invariance uses
\[
 \widetilde\alpha_g
  =\alpha'_g\,\frac{s_{1g}(\widetilde\gamma)^2}
                       {s_{2g}(\widetilde\gamma)}.
\]
It need not preserve a \emph{fixed original} $\alpha_g$; shape
variation can then carry additional information.

The singular curvature is also visible directly.
In \eqref{eq:effective-direct}, $D_gc_B=X^\top w$, so
$F_g=D_gc_Bl^\top$. At a spatial-intercept stationary point,
$0=c_B^\top U_\beta=\sum_j\widetilde s_j$ and hence
\[
 E_g=\left(\sum_jw_j\right)ll^\top
     =F_g^\top D_g^{-1}F_g.
\]
Summing over groups gives $\mathcal S=0$ for the covariate border
when each $D_g$ is invertible. Under the matched law the same
cancellation holds for expected information because $\E\widetilde s_j=0$.
Thus a positive-definite full regression inverse cannot be assumed
for this unrestricted totals-only setup.

This is a consequence of the observation model written in
paper Equation~9, not an assertion about an implementation that
also retains study-level allocation information.
The full two-margin Poisson and full clustered likelihoods retain
that information; the totals-only matched NB likelihood does not.
Scientifically specified restrictions or additional data can change
identifiability. A numerical ridge, or a roughness penalty that
leaves constant maps unpenalized, does not establish identification
from the original totals likelihood.

\subsection{Expected information under the matched distribution}
If the matched NB law is itself assumed, the outer scores have mean zero.
In particular $\E\Psi_{r'}=\E\Psi_A=0$, so inner-map curvature terms
drop \emph{in expectation}. At ${\mu_j}={r'} \mu_j^X/A$ the necessary expectations are
\begin{align}
 \E w_j&=({r'}+{\mu_j})\mu_j^XA/d_j^2={r'} \mu_j^X/d_j,\\
 \E\Psi_{{r'}{r'}}&=N\psi_1({r'})-\sum_j\E\{\psi_1(Y_j+{r'})\},\\
 \E\Psi_{{r'} A}&=-N/A+\sum_jd_j^{-1},\\
 \E\Psi_{AA}&=N{r'}/A^2-\sum_j({r'}+{\mu_j})/d_j^2,\\
 \E\Psi_{\eta_j{r'}}&=\mu_j^X/d_j,&
 \E\Psi_{\eta_j A}&=-({r'}+{\mu_j})\mu_j^X/d_j^2.
\end{align}
Insert these into \eqref{eq:aggregate-global}--\eqref{eq:aggregate-cross}.
The trigamma expectation is an NB-weighted sum over nonnegative counts;
it cannot generally be replaced by $\psi_1({\mu_j}+{r'})$.

Under the \emph{original} unequal-probability independent-cell sums,
the first two moments agree, but the expected digamma in the shape score
need not agree. Thus the matched shape score need not be centered under
the original distribution. Since ${r'}$ depends on $\gamma$ at fixed
original dispersion, the matched global regression score need not be
centered at the original parameter either.
An exact working Hessian does not resolve this possible change of target.
A robust variance, when justified, quantifies a working estimator's
variability; it does not automatically remove approximation bias.

\subsection{Stable log coordinates and exact recurrence checks}
Let $t_j=\eta_j-a$, ${p'}_j=(1+e^{-t_j})^{-1}$,
$\bar {p'}_j=1-{p'}_j$. Then \eqref{eq:Psi} becomes
\begin{equation}
 \Psi=-\sum_j\log\Gamma(Y_j+{r'})+N\log\Gamma({r'})
       +\sum_j({r'}+Y_j)\softplus(t_j)-\sum_jY_jt_j .
\end{equation}
Use $\softplus(t)=\max(t,0)+\log(1+e^{-|t|})$ for evaluation.
In coordinates $(\eta,q,a)$, direct differentiation gives
\begin{align}
 \Psi_a&=\sum_j\{Y_j-({r'}+Y_j){p'}_j\},\\
 \Psi_{aa}&=\sum_jw_j,\quad w_j=({r'}+Y_j){p'}_j\bar {p'}_j,\\
 \Psi_q&={r'}\left\{\sum_j\softplus(t_j)
                  -\sum_j[\psi(Y_j+{r'})-\psi({r'})]\right\},\\
 \Psi_{qa}&=-{r'}\sum_j{p'}_j,\\
 \Psi_{qq}&=\Psi_q+{r'}^2\sum_j\{\psi_1({r'})-\psi_1(Y_j+{r'})\},\\
 \Psi_{\eta_jq}&={r'} {p'}_j,\qquad\Psi_{\eta_ja}=-w_j .
\end{align}
In particular $\partial_q({r'}\Psi_{r'})=\Psi_q+{r'}^2\Psi_{{r'}{r'}}$:
the first derivative appears again through reparameterization.
An equivalent form of the regression blocks is
\begin{align}
 F_g&=X^\top({r'}p')d_q^\top-X^\top w\,d_a^\top,\\
 E_g&=\Psi_{qq}d_qd_q^\top+
       \Psi_{qa}(d_qd_a^\top+d_ad_q^\top)+\Psi_{aa}d_ad_a^\top
       +\Psi_qQ_q+\Psi_aQ_a . \label{eq:log-composition}
\end{align}
This is a rearrangement, not a smaller-term Hessian approximation.
For integer $y\geq0$, gamma recurrence yields exact finite sums
\begin{align}
 \log\Gamma({r'}+y)-\log\Gamma({r'})&=\sum_{k=0}^{y-1}\log({r'}+k),\\
 \psi({r'}+y)-\psi({r'})&=\sum_{k=0}^{y-1}({r'}+k)^{-1},\\
 \psi_1({r'})-\psi_1({r'}+y)&=\sum_{k=0}^{y-1}({r'}+k)^{-2}.
\end{align}
Empty sums are zero. These provide manual checks and can avoid subtracting
nearby special-function values. Log-sum-exp evaluation of
$\log s_1$ and $\log s_2$ likewise avoids overflowing the $\mu_i^Z$.

\subsection{Explicit optional log-dispersion blocks}
For a joint matched-model calculation let $\delta=\log\alpha$ in one group.
Then $q_\delta=a_\delta=-1$ and
$q_{\gamma\delta}=a_{\gamma\delta}=q_{\delta\delta}=a_{\delta\delta}=0$.
The complete additional negative-log-likelihood derivatives are
\begin{align}
 \Psi_\delta&=-\Psi_q-\Psi_a,\\
 H_{\delta\delta}&=\Psi_{qq}+2\Psi_{qa}+\Psi_{aa},\\
 H_{\gamma\delta}
 &=-(\Psi_{qq}+\Psi_{qa})d_q-(\Psi_{qa}+\Psi_{aa})d_a,\\
 H_{\beta\delta}&=-X^\top({r'}p')+X^\top w .
 \label{eq:aggregate-dispersion}
\end{align}
Different group dispersions have no direct cross derivatives.
The regression blocks remain \eqref{eq:log-composition}, including
their inner curvature terms. For effective log dispersion
$\delta'=\log\alpha'$, the same outer formulas hold but with the inner
derivatives \eqref{eq:effective-fixed}.
These formulas state both joint parameterizations without silently
changing what ``fixed dispersion'' means.

\section{Clustered NB: integrate the shared study multiplier}\label{sec:cluster}

\subsection{Gamma integral and full covariance}
For $i\in\mathcal I_g$, put $\nu=\nu_g=1/\alpha_g$ and assume
\begin{equation}
 \lambda_i\sim\operatorname{Gamma}(\text{shape }\nu,\text{rate }\nu),
 \qquad Y_{ij}\mid\lambda_i\ \text{independent Poisson}(\lambda_i\mu_{ij}).
\end{equation}
Different studies have independent multipliers.
Let $Y_{ti}=\sum_jY_{ij}$ and ${\mu_{ti}}=\sum_j\mu_{ij}=\mu_i^Zh_g$.
Multiply the conditional mass by the gamma density and integrate:
\begin{align}
 p(Y_i)&=\frac{\nu^\nu\prod_j\mu_{ij}^{Y_{ij}}}
                   {\Gamma(\nu)\prod_jY_{ij}!}
       \int_0^\infty \lambda^{Y_{ti}+\nu-1}e^{-({\mu_{ti}}+\nu)\lambda}\,d\lambda\\
 &=\frac{\Gamma(Y_{ti}+\nu)}{\Gamma(\nu)}
       \frac{\nu^\nu}{({\mu_{ti}}+\nu)^{Y_{ti}+\nu}}
       \prod_j\frac{\mu_{ij}^{Y_{ij}}}{Y_{ij}!}. \label{eq:cluster-mass}
\end{align}
The gamma integral is
$\int_0^\infty x^{a-1}e^{-bx}\,dx=\Gamma(a)b^{-a}$, $a,b>0$.
Its log is
\begin{equation}
 \ell_i=\log\Gamma(Y_{ti}+\nu)-\log\Gamma(\nu)+\nu\log\nu
 -(Y_{ti}+\nu)\log({\mu_{ti}}+\nu)+\sum_jY_{ij}\eta_{ij}
 -\sum_j\log(Y_{ij}!). \label{eq:cluster-ll}
\end{equation}
Summing \eqref{eq:cluster-ll} and using both data margins recovers
paper Equation~11 up to data-only factorials.

The law of total covariance gives, including diagonal entries,
\begin{align}
 \Cov(Y_{ij},Y_{ik})
 &=\E\{\Cov(Y_{ij},Y_{ik}\mid\lambda_i)\}
   +\Cov(\lambda_i\mu_{ij},\lambda_i\mu_{ik})\\
 &=\mathbf1_{\{j=k\}}\mu_{ij}+\alpha_g\mu_{ij}\mu_{ik},\\
 \Var(Y_i)&=\diag(\mu_i)+\alpha_g\mu_i\mu_i^\top . \label{eq:cluster-variance}
\end{align}
For distinct voxels this reduces to paper Equation~10.
Different studies have zero covariance under the model.
The diagonal is $\mu_{ij}+\alpha_g\mu_{ij}^2$, not just
$\alpha_g\mu_{ij}^2$. Also $\Var(Y_{ti})={\mu_{ti}}+\alpha_g{\mu_{ti}}^2$.
This construction captures shared study fluctuation; it does not introduce
an additional independent voxel-level overdispersion component.

\subsection{Exact factorization into a study total and spatial allocation}
Conditional on $\lambda_i$, the total has law
$Y_{ti}\mid\lambda_i\sim\operatorname{Poisson}(\lambda_i\mu_{ti})$.
The same gamma integral, now applied to this single count, gives
\[
 p(Y_{ti})=
 \frac{\Gamma(Y_{ti}+\nu)}{\Gamma(\nu)Y_{ti}!}
 \frac{\nu^\nu\mu_{ti}^{Y_{ti}}}{(\mu_{ti}+\nu)^{Y_{ti}+\nu}}.
\]
Thus $Y_{ti}$ is exactly NB with mean $\mu_{ti}$ and shape $\nu$.
For a fixed total the independent conditional Poisson counts have
multinomial probabilities
\[
 \pi_{ij}=\frac{\lambda_i\mu_{ij}}{\lambda_i\mu_{ti}}
         =\frac{\mu_{ij}}{\mu_{ti}}
         =\frac{\mu_{gj}^X}{h_g},\qquad g=g(i).
\]
The multiplier cancels, so integrating it does not change this
conditional multinomial distribution. Multiplying the two masses,
\begin{align}
 p(Y_i)
 &=p(Y_{ti})\,\frac{Y_{ti}!}{\prod_jY_{ij}!}
                     \prod_j\left(\frac{\mu_{ij}}{\mu_{ti}}\right)^{Y_{ij}}
 \notag\\
 &=\frac{\Gamma(Y_{ti}+\nu)}{\Gamma(\nu)}
   \frac{\nu^\nu}{(\mu_{ti}+\nu)^{Y_{ti}+\nu}}
   \prod_j\frac{\mu_{ij}^{Y_{ij}}}{Y_{ij}!}.
\end{align}
The factorial $Y_{ti}!$ and the power
$\mu_{ti}^{Y_{ti}-\sum_jY_{ij}}=1$ cancel explicitly, reproducing
\eqref{eq:cluster-mass}. Unlike the independent-cell sum in
Section~\ref{sec:aggregate}, this factorization uses no moment
approximation.

It also separates the sources of regression information.
Let $a_g=X^\top\mu_g^X$ as above. Then
\[
 \nabla_{\beta_g}\log\mu_{ti}=a_g/h_g,\qquad
 \nabla_\gamma\log\mu_{ti}=Z_i^\top.
\]
The NB total score on its log-mean scale is
$u_{ti}=(Y_{ti}-\mu_{ti})/(1+\alpha_g\mu_{ti})$.
The conditional allocation log likelihood, apart from data terms, is
$Y_i^\top X\beta_g-Y_{ti}\log h_g$.
Its spatial score is $X^\top Y_i-Y_{ti}a_g/h_g$, and its $\gamma$
score is zero. Combining total and allocation scores gives
\begin{align}
 U_{\beta_g,i}
 &=X^\top Y_i+(u_{ti}-Y_{ti})a_g/h_g
   =X^\top Y_i-c_i\mu_i^Z a_g,\\
 U_{\gamma,i}
 &=u_{ti}Z_i^\top=(Y_{ti}-c_i\mu_{ti})Z_i^\top,
 \qquad c_i=\frac{Y_{ti}+\nu_g}{\mu_{ti}+\nu_g}.
\end{align}
The cancellation follows from
$u_{ti}-Y_{ti}=-\mu_{ti}(Y_{ti}+\nu_g)/(\mu_{ti}+\nu_g)$.
The total captures the constant study effect, while conditional
allocation captures the normalized spatial shape.

\subsection{Study score and observed information on the log-mean scale}
At fixed $\nu$, define $c_i=(Y_{ti}+\nu)/({\mu_{ti}}+\nu)$.
Since $\partial {\mu_{ti}}/\partial\eta_{ij}=\mu_{ij}$,
\begin{equation}
 s_{ij}=\frac{\partial\ell_i}{\partial\eta_{ij}}
          =Y_{ij}-c_i\mu_{ij}.
\end{equation}
For any voxel $k$,
$\partial c_i/\partial\eta_{ik}=-c_i\mu_{ik}/({\mu_{ti}}+\nu)$.
The product rule therefore gives
\begin{align}
 -\frac{\partial s_{ij}}{\partial\eta_{ik}}
 &=c_i\mathbf1_{\{j=k\}}\mu_{ij}
              -\frac{c_i\mu_{ij}\mu_{ik}}{{\mu_{ti}}+\nu},\\
 W_i^{\rm obs}
 &=c_i\left\{\diag(\mu_i)-\frac{\mu_i\mu_i^\top}{{\mu_{ti}}+\nu}\right\}.
 \label{eq:cluster-W}
\end{align}
The second term is a rank-one correction, not a diagonal NB weight.
For $a\ne0$, weighted Cauchy--Schwarz gives
$(\sum_j\mu_{ij}a_j)^2\leq {\mu_{ti}}\sum_j\mu_{ij}a_j^2$, hence
\begin{equation}
 a^\top W_i^{\rm obs}a
 \geq c_i\frac{\nu}{{\mu_{ti}}+\nu}\sum_j\mu_{ij}a_j^2>0 .
\end{equation}
Thus the study log-mean curvature is positive definite at finite positive
means and $\nu>0$. Projection can still lose rank through the design.
Since $\E Y_{ti}={\mu_{ti}}$, $\E c_i=1$, giving
\begin{equation}
 W_i^{\rm F}=\diag(\mu_i)-\frac{\mu_i\mu_i^\top}{{\mu_{ti}}+\nu}.
\end{equation}
These weights and scores tend to the Poisson ones as $\nu\to\infty$.
The positive curvature establishes fixed-dispersion convexity only.

\subsection{Project to the paper's group and global coefficients}
Use $h_g$, $a_g=X^\top \mu_g^X$, and $B_g=X^\top\diag(\mu_g^X)X$
from the Poisson calculation.
The study's nonlinear negative-log term is
\begin{equation}
 f_i(t)=(Y_{ti}+\nu_g)\log(t+\nu_g),\quad
 f_i'(t)=c_i,\quad f_i''(t)=-\frac{c_i}{t+\nu_g}.
\end{equation}
The total ${\mu_{ti}}=\mu_i^Zh_g$ satisfies
\begin{align}
 \nabla_{\beta_g}{\mu_{ti}}&=\mu_i^Za_g,&
 \nabla_{\beta_g}^2{\mu_{ti}}&=\mu_i^ZB_g,\\
 \nabla_\gamma {\mu_{ti}}&={\mu_{ti}}{Z_i^\top},&
 \nabla_\gamma^2{\mu_{ti}}&={\mu_{ti}}{Z_i^\top}{Z_i},&
 \nabla_{\beta_g\gamma}^2{\mu_{ti}}&=\mu_i^Za_g{Z_i}.
\end{align}
The elementary composite rule is
$\nabla^2 f(t)=f'(t)\nabla^2t+f''(t)\nabla t\nabla t^\top$.
If $f$ is instead viewed as a function of $h_g$, its spatial block
is exactly $f_h' B_g+f_h''a_ga_g^\top$, with
$f_h'=c_i\mu_i^Z$ and $f_h''=-c_i(\mu_i^Z)^2/({\mu_{ti}}+\nu_g)$.
Thus the full study contributions are
\begin{align}
 D_{g,i}&=c_i\mu_i^ZB_g-\frac{c_i(\mu_i^Z)^2}{{\mu_{ti}}+\nu_g}a_ga_g^\top,\\
 F_{g,i}&=c_i\mu_i^Za_g{Z_i}-
       \frac{c_i\mu_i^Z{\mu_{ti}}}{{\mu_{ti}}+\nu_g}a_g{Z_i}
       =\frac{c_i\mu_i^Z\nu_g}{{\mu_{ti}}+\nu_g}a_g{Z_i},\\
 E_i&=c_i{\mu_{ti}}{Z_i^\top}{Z_i}-
       \frac{c_i{\mu_{ti}}^2}{{\mu_{ti}}+\nu_g}{Z_i^\top}{Z_i}
       =\frac{c_i{\mu_{ti}}\nu_g}{{\mu_{ti}}+\nu_g}{Z_i^\top}{Z_i} . \label{eq:cluster-blocks}
\end{align}
The scores, retaining the data terms, are
\begin{equation}
 U_{\beta_g,i}=X^\top Y_i-c_i\mu_i^Za_g,\qquad
 U_{\gamma,i}=(Y_{ti}-c_i{\mu_{ti}}){Z_i^\top} .
\end{equation}
Sum $D_{g,i},F_{g,i}$ over studies in group $g$ and $E_i$ over all studies.
Replace $c_i$ by one for expected information.
No spatially varying moderator design has been introduced.
The two data margins suffice for the total likelihood and derivatives
at fixed dispersion; a study-level robust calculation also needs
the individual study projections $X^\top Y_i$.

\subsection{Cluster shape and log-dispersion derivatives}
Differentiate \eqref{eq:cluster-ll} at fixed ${\mu_{ti}}$:
\begin{align}
 \ell_{i,\nu}&=\psi(Y_{ti}+\nu)-\psi(\nu)+\log\nu+1
                 -\log({\mu_{ti}}+\nu)-\frac{Y_{ti}+\nu}{{\mu_{ti}}+\nu},\\
 \ell_{i,\nu\nu}&=\psi_1(Y_{ti}+\nu)-\psi_1(\nu)+\frac1\nu
              -\frac1{{\mu_{ti}}+\nu}-\frac{{\mu_{ti}}-Y_{ti}}{({\mu_{ti}}+\nu)^2},\\
 \ell_{i,\eta_{ij}\nu}
 &=\frac{\mu_{ij}(Y_{ti}-{\mu_{ti}})}{({\mu_{ti}}+\nu)^2}.
\end{align}
The last derivative follows from differentiating $-c_i\mu_{ij}$.
For $\delta_g=\log\alpha_g$ and $\nu_g=e^{-\delta_g}$,
\begin{align}
 \ell_{\delta_g}&=-\nu_g\ell_{\nu_g},&
 \ell_{\delta_g\delta_g}&=\nu_g^2\ell_{\nu_g\nu_g}
                                      +\nu_g\ell_{\nu_g},\\
 \ell_{\vartheta\delta_g}&=-\nu_g\ell_{\vartheta\nu_g}. \label{eq:logdisp-chain}
\end{align}
Study derivatives are summed within the group before joint blocks are
formed. At a regular group shape optimum $\ell_{\nu_g}=0$,
but the extra score term in \eqref{eq:logdisp-chain} cannot be omitted
away from it. Expected shape/regression cross information is zero
for the correctly specified clustered model because $\E(Y_{ti}-{\mu_{ti}})=0$.
That fact need not hold for a misspecified distribution or observed Hessian.

\section{Optimization, fixed slices, and nuisance adjustment}\label{sec:optimization}

\subsection{Interpret Algorithm 1 as joint estimation by alternating updates}
The paper initializes the overdispersed fits from a Poisson solution
and alternates dispersion updates with regression updates using L-BFGS.
Mathematically, for a negative log likelihood $L(\vartheta,\delta)$,
one cycle takes the schematic form
\begin{equation}
 \delta^{(k+1)}\approx\arg\min_\delta L(\vartheta^{(k)},\delta),\qquad
 \vartheta^{(k+1)}\approx\arg\min_\vartheta L(\vartheta,\delta^{(k+1)}).
\end{equation}
Separate subproblems do not mean the final dispersion was externally known.
Check the full score, objective progress, and sensitivity to initialization.
Small objective changes alone can reflect a stalled or poorly scaled fit.
L-BFGS's limited-memory curvature approximation is an optimization device,
not automatically the observed-information covariance used for inference.

Poisson, underlying independent NB, and clustered NB have convex
negative log likelihood in the affine regression coefficients at fixed
dispersion, as proved above. This does not imply joint convexity in
regression and dispersion, nor convexity of the matched NB composition
at fixed original dispersion. It also does not guarantee numerical
convergence, full rank, or a finite estimate.
These distinctions refine, rather than replace, the paper's caution
about fitting its overdispersed models.

\subsection{Joint versus fixed-dispersion regression covariance}
Let $\delta=(\log\alpha_1,\ldots,\log\alpha_G)^\top$ and partition a
joint observed or expected information matrix in auxiliary coordinate
order $(\vartheta,\delta)$ as
\begin{equation}
 \mathcal H=\begin{pmatrix}
 H_{\vartheta\vartheta}&\mathcal B\\
 \mathcal B^\top&A_\delta
 \end{pmatrix}.
\end{equation}
This is a reordering and log transformation of the relevant full
paper parameter vector $\theta$, not a new meaning for that vector.
Use $\log\alpha'_g$ instead if effective dispersion is the coordinate
being held fixed, and form all blocks in that same coordinate system.
Assume the indicated blocks are invertible and the full matrix is
positive definite on the identified coordinates.
For a right-hand side $(b,0)$ the nuisance equation is
$\mathcal B^\top x_\vartheta+A_\delta x_\delta=0$, so
$x_\delta=-A_\delta^{-1}\mathcal B^\top x_\vartheta$.
Substitute into the regression equation:
\begin{equation}
 (H_{\vartheta\vartheta}-\mathcal B A_\delta^{-1}\mathcal B^\top)x_\vartheta=b.
\end{equation}
Therefore
\begin{equation}
 [\mathcal H^{-1}]_{\vartheta\vartheta}
   =(H_{\vartheta\vartheta}-\mathcal B A_\delta^{-1}\mathcal B^\top)^{-1},
 \qquad V_{\vartheta\mid\delta}\approx H_{\vartheta\vartheta}^{-1}.
 \label{eq:nuisance-schur}
\end{equation}
The second expression treats dispersion as fixed for inference.
The first allows local nuisance estimation uncertainty.
With positive-definite blocks it is no smaller in the PSD ordering.
Expected orthogonality can make $\mathcal B=0$ in particular correctly specified
models, but does not make an observed $\mathcal B$ identically zero.
The aggregate model needs its own derivatives, not a borrowed orthogonality.

For any smooth parameter transformation $\xi=\xi(\zeta)$,
\begin{equation}
 \nabla_\zeta^2L
 =J_\xi^\top(\nabla_\xi^2L)J_\xi+
       \sum_k L_{\xi_k}\nabla_\zeta^2\xi_k .
 \label{eq:general-transform}
\end{equation}
This explains both the log-dispersion extra score term and the
aggregate inner curvature. At $\alpha=0$, log dispersion is at infinity;
ordinary interior joint-Wald asymptotics do not address that boundary.
All inverse statements require the appropriate regular parameter space.

\section{Group blocks, Schur solves, and complete regression covariance}
\label{sec:covariance}

\subsection{Structural zeros do not imply independent fitted groups}
For each of the paper's likelihoods, with group-specific dispersion fixed,
the regression information has the structure
\begin{equation}
 H=\begin{pmatrix}D&F\\F^\top&E\end{pmatrix},\quad
 D=\blkdiag(D_1,\ldots,D_G),\quad
 F=\begin{pmatrix}F_1\\ \vdots\\ F_G\end{pmatrix}.
 \label{eq:border}
\end{equation}
Here $D_g$ is $P\times P$, $F_g$ is $P\times R$, and $E$ is $R\times R$.
The expressions are \eqref{eq:poisson-blocks},
\eqref{eq:aggregate-global}--\eqref{eq:aggregate-cross}, or
\eqref{eq:cluster-blocks}, for the chosen family.
Observed and expected versions must not be mixed without stating why.
Joint nuisance blocks can be retained in a larger border or removed
by the exact adjustment \eqref{eq:nuisance-schur}.
The unrestricted matched-totals setup in
Section~\ref{sec:aggregate-identification} has a singular covariate
border; the invertible formulas below cannot be applied to that case
without additional scientifically justified identifying restrictions.

When $R=0$, the inverse is $\blkdiag(D_g^{-1})$.
For $R>0$, common $\gamma$ connects groups through $F$.
The appropriate covariance is not obtained by dropping those connections.
Independently fitting each group's own $\gamma_g$ would instead define
a different model and, for independent datasets, independent fitted groups.
It is not the joint shared-covariate fit with a block-diagonal inverse.

\subsection{Derive elimination from the two block equations}
Solve $Hx=b$ with $x=(x_D,x_\gamma)$ and $b=(b_D,b_\gamma)$.
The equations are
\begin{align}
 Dx_D+Fx_\gamma&=b_D,\\
 F^\top x_D+Ex_\gamma&=b_\gamma .
\end{align}
Assume $D$ is nonsingular. Solve the first and substitute:
\begin{align}
 x_D&=D^{-1}b_D-D^{-1}Fx_\gamma,\\
 (E-F^\top D^{-1}F)x_\gamma&=b_\gamma-F^\top D^{-1}b_D.
\end{align}
Define $\mathcal T=D^{-1}F$ and the Schur complement
$\mathcal S=E-F^\top\mathcal T$.
The calligraphic $\mathcal T_g=D_g^{-1}F_g$ is a $P\times R$
matrix, distinct from the scalar $T_g=\sum_{i\in\mathcal I_g}\mu_i^Z$.
With ${\mathcal S}$ nonsingular, the algorithm is
\begin{equation}
 y=D^{-1}b_D,\qquad
 x_\gamma={\mathcal S}^{-1}(b_\gamma-F^\top y),\qquad x_D=y-\mathcal T x_\gamma.
 \label{eq:schur-solve}
\end{equation}
Each occurrence of an inverse multiplying a vector or matrix means a
linear solve with a factorization, not a demand to construct a dense inverse.
The block structure allows the group solves to be performed separately.

\subsection{Read off the inverse and prove definiteness conditions}
The coefficients of the right-hand sides in \eqref{eq:schur-solve} give
\begin{equation}
 H^{-1}=
 \begin{pmatrix}
 D^{-1}+\mathcal T\mathcal S^{-1}\mathcal T^\top&-\mathcal T\mathcal S^{-1}\\
 -\mathcal S^{-1}\mathcal T^\top&\mathcal S^{-1}
 \end{pmatrix}. \label{eq:schur-inverse}
\end{equation}
Multiplication verifies the identity. For example the upper-right block
of $HH^{-1}$ is $-D\mathcal T\mathcal S^{-1}+F\mathcal S^{-1}=0$,
since $D\mathcal T=F$.
The lower-right block is
$-F^\top\mathcal T\mathcal S^{-1}+E\mathcal S^{-1}
=\mathcal S\mathcal S^{-1}=I_R$.
The remaining blocks similarly give $I_{GP}$ and zero.

For vectors $a,b$, complete the square:
\begin{equation}
 \begin{pmatrix}a\\b\end{pmatrix}^{\!\top}H
 \begin{pmatrix}a\\b\end{pmatrix}
 =(a+\mathcal T b)^\top D(a+\mathcal T b)+b^\top\mathcal S b.
 \label{eq:complete-square}
\end{equation}
Thus with $D\succ0$, $H\succ0$ if and only if ${\mathcal S}\succ0$.
An invertible but indefinite matrix also has an algebraic inverse,
but that inverse is not a valid likelihood covariance.
For regular likelihood inference the covariance approximation is
$V\approx H(\widehat\vartheta)^{-1}$ or $I(\widehat\vartheta)^{-1}$,
subject to the model, asymptotic regime, and nuisance convention.

\subsection{Shared covariates induce a low-rank group adjustment}
Write $\mathcal T_g=D_g^{-1}F_g$. Equation \eqref{eq:schur-inverse} gives
\begin{align}
 V_{\beta_g\beta_h}
 &=\mathbf1_{\{g=h\}}D_g^{-1}+\mathcal T_g\mathcal S^{-1}\mathcal T_h^\top,\\
 V_{\beta_g\gamma}&=-\mathcal T_g\mathcal S^{-1},&
 V_{\gamma\gamma}&={\mathcal S}^{-1}. \label{eq:cov-blocks}
\end{align}
The spatial adjustment $\mathcal T\mathcal S^{-1}\mathcal T^\top$
is PSD and has rank at most $R$
when ${\mathcal S}\succ0$. Its individual cross-group entries need not be positive.
For a two-group spatial coefficient contrast,
\begin{equation}
 \Var(\widehat\beta_g-\widehat\beta_h)
 \approx D_g^{-1}+D_h^{-1}
       +(\mathcal T_g-\mathcal T_h)\mathcal S^{-1}
        (\mathcal T_g-\mathcal T_h)^\top .
\end{equation}
Using the two marginal variances but dropping their covariance is wrong;
using only $D_g^{-1}+D_h^{-1}$ omits shared-nuisance uncertainty unless
the adjustment cancels. This cancellation can occur for a contrast even
when each individual map has substantial nuisance-induced uncertainty.

\subsection{Selected contrasts and marginal diagonals without dense inversion}
For an auxiliary $k\times(GP+R)$ coefficient contrast $\mathcal C$, solve
$HB=\mathcal C^\top$ by \eqref{eq:schur-solve}. Then
\begin{equation}
 \mathcal C V\mathcal C^\top\approx\mathcal C B.
 \label{eq:selected-solve}
\end{equation}
For one contrast column $c$, solve $Hw_c=c$ and use $c^\top w_c$.
For selected coefficient variances set $c$ to the corresponding unit
vectors; for map variances set it to the spline-evaluation vectors.
One may batch right-hand sides in memory-limited chunks.

Partition $\mathcal C=(\mathcal C_D,\mathcal C_\gamma)$ and insert
\eqref{eq:schur-inverse}. These are coefficient-space blocks, not
the paper's selected-group matrix $C$ or selected-covariate matrix
$C_\gamma$:
\begin{equation}
 \mathcal C H^{-1}\mathcal C^\top
 =\mathcal C_DD^{-1}\mathcal C_D^\top+
 (\mathcal C_\gamma-\mathcal C_D\mathcal T)\mathcal S^{-1}
 (\mathcal C_\gamma-\mathcal C_D\mathcal T)^\top .
 \label{eq:contrast-factor}
\end{equation}
This directly displays nuisance cancellation.
For group $g$ at voxel $j$,
\begin{equation}
 \Var(x_j^\top\widehat\beta_g)
 \approx x_j^\top D_g^{-1}x_j+
        (\mathcal T_g^\top x_j)^\top\mathcal S^{-1}(\mathcal T_g^\top x_j).
\end{equation}
If $D_g=L_gL_g^\top$ is a Cholesky factorization, its first term is
$\|L_g^{-1}x_j\|^2$. Analogous triangular solves evaluate the second.
An $N\times N$ map covariance is not needed for $N$ standard errors.
No arbitrary eigenvalue clipping or diagonal perturbation is hidden in
these formulas.

\subsection{Why the Schur complement is profile curvature}
For a smooth local spatial optimum $\beta_*(\gamma)$,
$L_\beta(\beta_*(\gamma),\gamma)=0$.
Differentiate this equation with respect to $\gamma$:
\[
 D\,\frac{\partial\beta_*}{\partial\gamma^\top}+F=0,
 \qquad \frac{\partial\beta_*}{\partial\gamma^\top}=-D^{-1}F.
\]
For the profiled negative log likelihood
$L_{\rm prof}(\gamma)=L(\beta_*(\gamma),\gamma)$, the chain rule gives
$\nabla_\gamma L_{\rm prof}=L_\gamma$ because $L_\beta=0$.
Differentiating once more yields
\[
 \nabla_\gamma^2 L_{\rm prof}
 =E+F^\top\frac{\partial\beta_*}{\partial\gamma^\top}
 =E-F^\top D^{-1}F=\mathcal S.
\]
This identity concerns the observed Hessian along a differentiable
interior solution branch with nonsingular $D$.
Using expected information gives the corresponding efficient-information
construction, not an exact data-specific profile derivative.

\subsection{A closed-form Poisson nuisance adjustment}
For the unpenalized full Poisson model, suppose the identified spatial
basis contains a constant map: $Xc_B=\mathbf1_N$.
Keep the spatial intercept available in each group and exclude
confounded constant study columns so the full model is identifiable.
Recall the scalar $T_g$, vector $t_g^{(z)}$, and matrix $V_g$
defined in Section~\ref{sec:poisson}. Then
\[
 B_gc_B=X^\top\diag(\mu_g^X)Xc_B
       =X^\top\mu_g^X=a_g,\qquad a_g^\top c_B=h_g.
\]
Consequently, whenever $D_g=T_gB_g$ is invertible,
\begin{align}
 \mathcal T_g
 &=D_g^{-1}F_g
   =\frac{c_B(t_g^{(z)})^\top}{T_g},\\
 F_g^\top D_g^{-1}F_g
 &=\frac{h_g}{T_g}t_g^{(z)}(t_g^{(z)})^\top,\\
 \mathcal S
 &=\sum_g h_g\left\{V_g-
       \frac{t_g^{(z)}(t_g^{(z)})^\top}{T_g}\right\}
   =\sum_g h_gT_g\Cov_{\pi_g}(Z_i^\top),
 \qquad \pi_{ig}=\mu_i^Z/T_g .
 \label{eq:poisson-efficient}
\end{align}
The last equality follows by writing the weighted covariance as
$\sum_i\pi_{ig}Z_i^\top Z_i-\bar z_g\bar z_g^\top$ with
$\bar z_g=t_g^{(z)}/T_g$.
Thus information on shared covariates after estimating group baselines
is exactly weighted \emph{within-group} variation in study covariates.
For any direction $a$,
\[
 a^\top\mathcal S a
  =\sum_g h_gT_g\sum_{i\in\mathcal I_g}
      \pi_{ig}\{Z_i a-\bar z_g^\top a\}^2.
\]
It vanishes precisely when $Z_i a$ is constant within each group.
This makes the identification condition explicit and computes the
small Schur matrix without large spatial solves under the stated
Poisson assumptions.
It is not a replacement for the general NB, penalized, or constrained
formulas.
At an interior Poisson optimum, the spatial intercept score additionally
gives $h_gT_g=Y_{g\cdot}$; then \eqref{eq:poisson-efficient} agrees
with the allocation information derived after Equation~\eqref{eq:allocation}.

\section{Roughness penalties and uncertainty interpretations}\label{sec:penalty}

\subsection{A declared quadratic convention, not an invented operator}
Paper Section 2.1.1 introduces a $P\times P$ roughness matrix $J$ and
refers its construction to Supplementary Appendix A.1.
That supplement is not supplied within the reference PDF used here.
We therefore do not assert an exact derivative operator, integration
domain, boundary condition, or normalization for its $J$.
The following generic calculation applies to a supplied fixed symmetric
$J\succeq0$ and declared nonnegative smoothing strengths $\tau_g$:
\begin{equation}
 \mathcal P(\vartheta)=\frac12\sum_g\tau_g\beta_g^\top J\beta_g
 =\frac12\vartheta^\top K\vartheta,\quad
 K=\blkdiag(\tau_1J,\ldots,\tau_GJ,0_{R\times R}).
 \label{eq:penalty}
\end{equation}
This explicitly chooses the factor $1/2$.
If a convention omits that factor, the same numerical $\tau_g$
produces twice the gradient and twice the curvature.

Expand one term as $\tau_g\sum_{ab}\beta_{ga}J_{ab}\beta_{gb}/2$.
Differentiating with respect to $\beta_{gc}$ gives
\begin{align}
 \partial_{\beta_{gc}}\mathcal P
 &=\frac{\tau_g}{2}
       \left(\sum_bJ_{cb}\beta_{gb}+\sum_a\beta_{ga}J_{ac}\right)
 =\tau_g(J\beta_g)_c,\\
 \partial_{\beta_{gc}\beta_{ge}}^2\mathcal P&=\tau_gJ_{ce}.
\end{align}
Thus for $L_{\rm pen}=-\ell+\mathcal P$,
\begin{equation}
 \nabla L_{\rm pen}=-U+K\vartheta,\qquad H_{\rm pen}=H+K.
\end{equation}
Separate group penalties preserve the group blocks and shared border.
The penalized score equation is $U(\widehat\vartheta)=K\widehat\vartheta$,
not $U(\widehat\vartheta)=0$.

\subsection{Sampling covariance versus posterior inverse curvature}
Expand the penalized equation near an unpenalized target $\vartheta_0$:
\begin{align}
 0&\approx U(\vartheta_0)-H(\vartheta_0)(\widehat\vartheta-\vartheta_0)
                 -K\vartheta_0-K(\widehat\vartheta-\vartheta_0),\\
 (H+K)(\widehat\vartheta-\vartheta_0)&\approx U(\vartheta_0)-K\vartheta_0 .
\end{align}
Under a correctly specified likelihood, replacing sensitivity by $I$
gives the local approximations
\begin{equation}
 \operatorname{bias}(\widehat\vartheta)\approx-(I+K)^{-1}K\vartheta_0,\qquad
 V_{\rm freq}\approx(I+K)^{-1}I(I+K)^{-1}.
 \label{eq:penalty-frequentist}
\end{equation}
Strong smoothing can change the population target; this linearization
is not a global exact bias identity.
By contrast a Gaussian-prior interpretation of $K$ leads to a Laplace
posterior covariance approximation $(H+K)^{-1}$.
Only an actually quadratic posterior would make this exact.
In expected-curvature notation the difference is
\begin{equation}
 (I+K)^{-1}-(I+K)^{-1}I(I+K)^{-1}
 =(I+K)^{-1}K(I+K)^{-1}\succeq0.
\end{equation}
Thus inverse penalized curvature alone is not generally the frequentist
sampling variance of a penalized estimator. Neither variance includes
squared smoothing bias about the unpenalized target.
For scalar information $i>0$ and penalty $k$, the two variances are
$1/(i+k)$ and $i/(i+k)^2$, respectively.

\subsection{Smoothing selection and identifiability remain separate}
All preceding derivatives condition on fixed $\tau_g$.
If a data-driven rule selects them, differentiating the score equation
along a smooth solution branch gives
\begin{equation}
 (H+K)\frac{\partial\widehat\vartheta}{\partial\tau_g}
 =-\frac{\partial K}{\partial\tau_g}\widehat\vartheta .
\end{equation}
A joint delta calculation must include uncertainty in $\widehat\tau$
and its covariance with the regression fit, or a bootstrap must repeat
the selection rule. Merely adding an independent variance term is
unjustified in general.
For $H,K\succeq0$, $H+K\succ0$ precisely when
$\ker(H)\cap\ker(K)=\{0\}$, because both quadratic forms are nonnegative.
A roughness penalty leaving constant maps unpenalized cannot resolve
the constant-map confounding direction in \eqref{eq:confounding}.
Adding a numerical ridge changes the inverse and possibly the estimator;
it is not evidence that the original information supported reliable tests.

\section{Group maps, covariate contrasts, and Wald inference}\label{sec:inference}

\subsection{Linear map evaluation and full voxelwise group covariance}
The baseline log-intensity map is $\eta_g^X=X\beta_g$, an $N$-vector.
For its coefficient covariance $V_{gg}$, the dimensionally correct identity is
\begin{equation}
 \Var(\widehat\eta_g^X)\approx XV_{gg}X^\top ,
 \label{eq:map-variance}
\end{equation}
not $X^\top V_{gg}X$ when $X$ is $N\times P$.
For all groups at voxel $j$, define the auxiliary evaluation matrix
\begin{equation}
 A_j=(I_G\otimes x_j^\top,\ 0_{G\times R}).
\end{equation}
Let $g_1,\ldots,g_S$ be the $S$ groups involved in the test and let
$P_G\in\R^{S\times G}$ select them in that order.
On the log-intensity scale recommended by the paper,
\begin{equation}
 \theta_j=
 \begin{pmatrix}\log\mu_{g_1j}^X\\ \vdots\\\log\mu_{g_Sj}^X\end{pmatrix}
 =P_GA_j\vartheta,\qquad
 V_j=P_GA_j V A_j^\top P_G^\top .
 \label{eq:voxel-selection}
\end{equation}
Here $\theta_j$ and $V_j$ now have exactly their paper Equation~12
meanings and dimensions. $V$ is the chosen regression covariance
estimate, so $V_j$ is its induced voxel covariance estimate.
Its $(a,b)$ entry is
$x_j^\top V_{\beta_{g_a}\beta_{g_b}}x_j$, including the cross-group
terms induced by estimating shared $\gamma$.
For the paper's contrast $C\in\R^{m\times S}$, define a new
\emph{coefficient-space} matrix $\mathcal C_j=CP_GA_j$.
Then
\begin{equation}
 C\theta_j=\mathcal C_j\vartheta,\qquad
 C V_j C^\top=\mathcal C_j V\mathcal C_j^\top .
 \label{eq:group-contrast}
\end{equation}
The covariance can therefore be computed by selected coefficient
solves without forming $V$ or any $N\times N$ map covariance.
This evaluates the spatial baseline at the zero covariate row.
For a prediction at fixed covariate row $Z_*\in\R^{1\times R}$,
append $Z_*$ to the appropriate spatial evaluation row.
For group contrasts whose row sums are zero, the common $Z_*\gamma$
term cancels; it does not cancel for arbitrary contrasts.

\subsection{Wald rank, degrees of freedom, and estimability}
For the paper's null $H_0:C\theta_j=0$, put
$d_j=C\widehat\theta_j$ and $\Omega_j=C V_j C^\top$.
With $m$ linearly independent estimable rows and $\Omega_j\succ0$,
\begin{equation}
 Q_j=(C\widehat\theta_j)^\top(C V_j C^\top)^{-1}
          (C\widehat\theta_j)
       \ \overset{H_0}{\approx}\ \chi_m^2 .
 \label{eq:wald}
\end{equation}
This is paper Equation~12; $Q_j$ is only a helper name for its
squared statistic. To see the degrees of freedom, take
$\Omega_j=L_jL_j^\top$. Under the null, asymptotic normality gives
$L_j^{-1}d_j\approx N_m(0,I_m)$, and
$Q_j=\|L_j^{-1}d_j\|^2$ is approximately a sum of $m$ independent
squared standard normals.
For a nonzero specified null $C\theta_j=d_0$, replace $d_j$
by $C\widehat\theta_j-d_0$.
The approximation requires asymptotic normality and an appropriate
covariance estimate, not merely an invertible numerical matrix.
Redundant rows must be removed, or a generalized-inverse version uses
the rank of the estimable covariance on its supported null space.
Unidentified directions cannot be made testable by a pseudoinverse alone.
With a singular design, an estimable contrast must annihilate its null
directions under the chosen constraints.

For one scalar contrast with positive standard error,
\begin{equation}
 W_j=\frac{C\widehat\theta_j}{\sqrt{C V_j C^\top}}
      \ \overset{H_0}{\approx}\ N(0,1),\qquad Q_j=W_j^2,
 \quad C\in\R^{1\times S}.
\end{equation}
This preserves the signed $W_j$ of paper Equation~13, instead of
renaming it $Z$ and colliding with the study design.
For example, take $S=2$, $C=(1,-1)$, and groups $g,h$.
Then
\[
 W_j=
 \frac{x_j^\top(\widehat\beta_g-\widehat\beta_h)}
 {\sqrt{x_j^\top
 (V_{\beta_g\beta_g}+V_{\beta_h\beta_h}
     -V_{\beta_g\beta_h}-V_{\beta_h\beta_g})x_j}}.
\]
The two cross-covariances follow directly by expanding the variance
of a difference, rather than assuming independent group estimates.
No universal ``sum-to-one'' condition is imposed on a contrast.
Differences such as $(1,-1)$ sum to zero, averages can sum to one,
and coefficient tests use other valid rows. The null determines the rows;
independence and estimability determine the degrees of freedom.

\subsection{Shared global-covariate inference}
Select $s$ of the $R$ global covariates with a row-selection matrix
$P_\gamma\in\R^{s\times R}$ and define
$\gamma_{\rm sel}=P_\gamma\gamma$.
The paper's $C_\gamma$ has dimension $m\times s$, not
$m\times R$ unless all covariates are selected.
The selection and contrast steps give
\[
 V_{\rm sel}=P_\gamma V_{\gamma\gamma}P_\gamma^\top,\qquad
 \Var(C_\gamma\widehat\gamma_{\rm sel})
       \approx C_\gamma V_{\rm sel} C_\gamma^\top .
\]
For $H_0:C_\gamma\gamma_{\rm sel}=0$, paper Equation~14 is therefore
\begin{equation}
 Q_\gamma=(C_\gamma\widehat\gamma_{\rm sel})^\top
       (C_\gamma V_{\rm sel}C_\gamma^\top)^{-1}
       (C_\gamma\widehat\gamma_{\rm sel})\ \approx\ \chi_m^2 .
\end{equation}
Under the fixed-dispersion inverse-information convention,
$V_{\rm sel}=P_\gamma\mathcal S^{-1}P_\gamma^\top$.
Joint dispersion adjustment, robust covariance, or penalized sampling
covariance can require the corresponding replacement.
All $m$ rows must again be independent and estimable.
For a specified nonzero null, subtract its $m$-vector on both sides
of the quadratic form.
The conditional inverse $E^{-1}$ is not generally $V_{\gamma\gamma}$.
Even a small $R$ does not guarantee stable covariate inference:
${\mathcal S}$ can be nearly singular when covariates are confounded with group
baselines or with one another. Nonzero fitted coefficients do not prove
information rank. Low dimension helps computation, not identification.
Because covariates have units and preprocessing scales, comparisons
between two components of $\gamma$ must specify that scale.

\subsection{Exponentiation and optional normalization}
For the baseline intensity $\mu_g^X=\exp(X\beta_g)$, componentwise differentiation
gives $d\mu_g^X=\diag(\mu_g^X)X\,d\beta_g$.
Therefore the first-order delta covariance is
\begin{equation}
 \Cov(\widehat \mu_g^X,\widehat \mu_h^X)
 \approx \diag(\mu_g^X)X V_{\beta_g\beta_h}X^\top\diag(\mu_h^X).
 \label{eq:exp-delta}
\end{equation}
At a voxel, $\Var(\widehat \mu_{gj}^X)\approx
(\mu_{gj}^X)^2\Var(x_j^\top\widehat\beta_g)$.
The map from $\beta_g$ to log intensity is linear; the exponentiation step
introduces an additional delta approximation.
If the paper's $\theta_j$ is instead taken on the intensity scale,
replace \eqref{eq:voxel-selection} by
$\theta_j=\exp(P_GA_j\vartheta)$ and use Jacobian
$\diag(\theta_j)P_GA_j$ to form $V_j$.
Using the log-map covariance unchanged in an intensity-scale test
would omit these mean factors.
A log-scale interval exponentiates to a positive intensity interval,
and a log-difference interval exponentiates to an intensity-ratio interval.

As an optional derived output, normalize $q_g=\mu_g^X/h_g$.
Since $dh_g=(\mu_g^X)^\top X\,d\beta_g$, the quotient rule gives
\begin{equation}
 \frac{\partial q_g}{\partial\beta_g^\top}
 =\{\diag(q_g)-q_gq_g^\top\}X .
\end{equation}
The row derivatives sum to zero because $\sum_jq_{gj}=1$.
Its covariance is necessarily singular in that normalization direction.
This optional probability map is not the same estimand as the paper's
unnormalized baseline intensity; normalize only if that is the intended
scientific question.

\subsection{Homogeneous references estimated from the data}
For a genuinely known homogeneous log baseline $\eta_{g0}$,
the scalar null $x_j^\top\beta_g=\eta_{g0}$ uses the standard contrast
variance. If $\eta_{g0}$ is estimated from the same data, instead
\begin{equation}
 \Var(x_j^\top\widehat\beta_g-\widehat\eta_{g0})
 =\Var(x_j^\top\widehat\beta_g)+\Var(\widehat\eta_{g0})
          -2\Cov(x_j^\top\widehat\beta_g,\widehat\eta_{g0}).
\end{equation}
The estimated reference is not automatically a fixed null constant.
For example, a linear reference $a^\top\beta_g$ is tested with
row $(x_j-a)^\top$, not with $x_j^\top$ alone.
A separately constrained homogeneous fit may require joint influence
calculations or a null bootstrap repeating the estimation.
For a concrete alternative reference derived from the fitted map,
take $\eta_{g0}=\log(h_g/N)$, the log of its mean intensity.
This is not asserted to be the same as a separately refitted null.
Because $dh_g=a_g^\top d\beta_g$,
\[
 \nabla_{\beta_g}\{x_j^\top\beta_g-\log(h_g/N)\}
      =x_j-a_g/h_g,
\]
so its local delta variance is
\[
 (x_j-a_g/h_g)^\top V_{\beta_g\beta_g}(x_j-a_g/h_g),
\]
evaluated at the fitted parameters. The centered log intensity is
$\log(Nq_{gj})$, which also exhibits its relation to the optional
normalized map.
The same qualification applies after exponentiating the reference.
Finally, voxelwise $p$-values are not already familywise-corrected results;
multiple-testing procedures require their own assumptions and a
clearly specified testing family.

\section{Resampling and an optional robust-covariance extension}\label{sec:bootstrap}

\subsection{Study sandwich: an extension, not the core paper calculation}
For a genuine study-factorized likelihood, let
$U_i=\nabla_\vartheta\ell_i$ and $H=-\sum_i\nabla_\vartheta U_i^\top$.
A standard study-cluster sandwich estimate is
\begin{equation}
 \widehat V_{\rm rob}=H^{-1}\left(\sum_iU_iU_i^\top\right)H^{-1}.
 \label{eq:sandwich}
\end{equation}
These are sums, so no additional $1/M$ factor is inserted.
Its usual justification uses independent studies, enough independent units,
identification, and appropriate estimating-equation regularity.
Poisson uses $U_i=\sum_jb_{ij}(Y_{ij}-\mu_{ij})$.
The underlying independent NB uses its weighted residual score, whereas
clustered NB uses $U_i=B_i^\top(Y_i-c_i\mu_i)$.
Substituting Poisson residuals into a fitted NB sandwich changes the
estimating equation and is not a general NB robust covariance.

The aggregate NB likelihood is a function of group--voxel totals and
nonlinear moment maps across studies. It does not provide arbitrary
study-factorized likelihood scores. A valid influence calculation for
that aggregate estimator, or refitting after study resampling, requires
separate justification. Assigning each study a convenient fraction of
an aggregate score does not establish \eqref{eq:sandwich}.

For selected coefficient contrasts solve $HB=\mathcal C^\top$ and accumulate
\begin{equation}
 \mathcal C\widehat V_{\rm rob}\mathcal C^\top
 =\sum_i(B^\top U_i)(B^\top U_i)^\top .
\end{equation}
This is a PSD Gram matrix with rank at most $\min(M,k,GP+R)$.
At an exactly solved unpenalized score equation, $\sum_iU_i=0$
reduces the bound to $M-1$. Few studies can therefore limit inference
even with many observed foci.
Joint nuisance estimation requires corresponding joint scores and bread.
For fixed penalties the bread is $H+K$, while the deterministic penalty
is not an extra independent random score in the meat.
Score centering and the penalized population target require attention
under substantial shrinkage; a naive uncentered meat is not universal.

\subsection{Null resampling tailored to the tested hypothesis}
Paper Section 2.1.4 describes randomization under spatial homogeneity
using a binomial process and group comparison using a pooled shared
intensity fit followed by regeneration and refitting.
These procedures simulate a specified null; they are not distribution-free.
A careful operational specification is:
\begin{enumerate}
 \item State the null, including which group maps coincide or are
 homogeneous, whether study totals are conditioned on, the voxel mask,
 covariates, dispersion convention, and the count-generating distribution.
 \item Fit the null, estimating its nuisance parameters under its
 constraints. For shared-baseline group comparisons retain the intended
 covariate design and the numbers of studies in each group.
 \item Generate independent studies under that null. For clustered NB,
 draw one gamma multiplier per study and then its conditionally
 independent counts; drawing one new multiplier per voxel changes the model.
 \item Apply the same preprocessing conventions and refit the same
 alternative or statistic-producing procedure used on the observed data.
 Re-estimate fitted dispersion, and repeat penalty selection if it was
 part of the original procedure rather than a prespecified constant.
 \item Compare the observed statistic with the simulated null statistics,
 using the chosen one-sided, two-sided, or maximum-statistic rule.
\end{enumerate}
For aggregate NB, distinguish simulating the matched totals distribution
from simulating original independent NB studies and then aggregating.
They assess different distributions except in an exact closure case.
If the statistic requires study totals or study-level resampling, aggregate
totals alone do not uniquely specify a study allocation.

\subsection{Replicates, failed fits, and what the foci threshold means}
For $B$ successful null replicates with statistics $T_b$, a common
Monte Carlo tail estimate is
\begin{equation}
 \widehat p=\frac{1+\sum_{b=1}^B\mathbf1_{\{T_b\geq T_{\rm obs}\}}}{B+1}.
\end{equation}
The comparison statistic must encode sidedness consistently.
The smallest attainable value is $1/(B+1)$, which matters for
multiple-testing resolution. Conditional on an ideal fixed null,
Monte Carlo variation is on the order of
$\sqrt{p(1-p)/B}$; nuisance estimation adds other approximation error.
The paper describes at least 1,000 repetitions.
This is a reported procedural choice, not a universal guarantee of
accurate extreme tails.

Predefine convergence checks and a failed-fit policy.
Record attempted and successful fits, retry rules, boundary estimates,
and sensitivity to failures. Silently deleting difficult fits can select
a nonrepresentative null distribution.
Resampling whole independent studies, typically within groups for
group-specific sampling targets, preserves their within-study dependence.
Voxelwise resampling does not.
A nonparametric study bootstrap and a parametric null bootstrap address
different questions; a hypothesis test needs an appropriate null mechanism.

The paper's practical criterion of approximately 200 foci for reliable
inverse-information inference is empirical guidance from its setting.
It is not a theorem about rank, coverage, or every design.
Rank and conditioning also depend on spline support, study allocation,
covariate variation, dependence, and fitted intensity.
Bootstrap avoids directly using some unstable inverses but cannot
automatically cure nonidentification, poor model fit, or failed optimization.

\section{Numerical audit, conditioning, and computational accounting}
\label{sec:numerics}

\subsection{Symmetric eigenvalues and scale-aware diagnostics}
For a symmetric positive-definite information matrix, the spectral
condition number is
\begin{equation}
 \kappa_2(H)=\lambda_{\max}(H)/\lambda_{\min}(H).
\end{equation}
To derive this without a general singular-value decomposition, let
$H=O\Lambda O^\top$ with $O$ orthogonal. Then
$H^\top H=O\Lambda^2O^\top$, so the singular values are
$|\lambda_a|$. For any symmetric nonsingular $H$,
\[
 \kappa_2(H)=\frac{\max_a|\lambda_a|}{\min_a|\lambda_a|};
\]
positive definiteness permits dropping the absolute values.
A symmetric eigensolver therefore targets the same spectral condition
number as a general SVD; it does not merely estimate a different norm.
Inspect signed symmetric eigenvalues, not just magnitudes: a negative
eigenvalue indicates indefinite curvature, whereas a tiny positive one
indicates weak information or scaling. Check symmetry before applying a
symmetric eigensolver; substantial antisymmetry indicates an assembly or
differentiation problem.
No universal threshold on $\kappa_2$ proves inferential validity.

Block conditioning alone is insufficient.
For $H=\diag(1,10^{-8})$, each scalar block has condition number one,
but the full matrix has condition number $10^8$.
For independent blocks the characteristic polynomial is the product
of their characteristic polynomials, so the full spectrum is their
union. Use the global largest and smallest eigenvalue magnitudes.
For a bordered matrix, Schur elimination is a congruence rather than
an orthogonal similarity: the spectra of $D$ and ${\mathcal S}$ do not generally
combine into the spectrum of $H$.
Border coupling can create a second problem:
\begin{equation}
 H_\varepsilon=\begin{pmatrix}1&1-\varepsilon\\1-\varepsilon&1\end{pmatrix},
 \quad \lambda(H_\varepsilon)=\{\varepsilon,2-\varepsilon\},\quad
 {\mathcal S}=1-(1-\varepsilon)^2 .
\end{equation}
For small positive $\varepsilon$, all diagonal blocks are benign,
yet the full covariance is unstable. This is near confounding through
the shared border. Examine $D_g$, ${\mathcal S}$, and scale, not just blockwise rank.

If $\vartheta=A\phi$ for a fixed nonsingular scale matrix $A$,
$H_\phi=A^\top H_\vartheta A$ and $V_\vartheta=AV_\phi A^\top$.
Scaling can improve numerical solves while leaving correctly transformed
contrasts unchanged. It cannot create information in a true null direction.
For a computed solve $Hx=b$, inspect the relative residual
$\|Hx-b\|/(\|H\|\|x\|+\|b\|)$; a small residual is a numerical check,
not a proof that estimated standard errors have correct coverage.

\subsection{Manual Schur example with two groups and one covariate}
Take $G=2$, $P=R=1$, with
\begin{equation}
 D=\begin{pmatrix}2&0\\0&3\end{pmatrix},\quad
 F=\begin{pmatrix}1\\1\end{pmatrix},\quad E=2.
\end{equation}
Then $\mathcal T=(1/2,1/3)^\top$, ${\mathcal S}=2-1/2-1/3=7/6$, and
\begin{equation}
 H=\begin{pmatrix}2&0&1\\0&3&1\\1&1&2\end{pmatrix},
 \qquad H^{-1}=\frac17
 \begin{pmatrix}5&1&-3\\1&3&-2\\-3&-2&6\end{pmatrix}.
\end{equation}
Both $D$ and ${\mathcal S}$ are positive definite, so \eqref{eq:complete-square}
proves $H\succ0$. Multiplying the displayed matrices gives identity.
The fitted groups have covariance $1/7$, despite their zero direct
information cross block. Their difference has variance
$(5+3-2)/7=6/7$.
Treating $\gamma$ as fixed instead gives $1/2+1/3=5/6$.
The shared-nuisance adjustment is
$(1/2-1/3)^2(6/7)=1/42$, so $5/6+1/42=6/7$.
Adding the marginal joint variances while omitting covariance would give
$8/7$, a different and incorrect result for this contrast.
These are audit numbers, not a fitted example or a performance claim.

\subsection{Manual moment-matching examples}
For two studies take $\mu_1^Z=1$, $\mu_2^Z=2$, $\mu_j^X=1/2$, and $\alpha=1$.
Then $s_1=3$, $s_2=5$, $s_3=9$, and
\begin{equation}
 {\mu_j}=\frac32,\quad \Var(Y_j)=\frac32+\frac54=\frac{11}{4},\quad
 {r'}=\frac95,\quad A=\frac35,\quad {p'}_j=\frac5{11}.
\end{equation}
The matched variance is
${\mu_j}+{\mu_j}^2/{r'}=3/2+(9/4)/(9/5)=11/4$,
confirming the two-moment calculation.
The third-cumulant discrepancy from \eqref{eq:third-cumulant} is
$2(1/2)^3(9-25/3)=1/6$, demonstrating nonexact closure.

For an exact common-probability case let $\mu_1^Z=\mu_2^Z=1$ with the same
$\mu_j^X=1/2$ and $\alpha=1$.
Then $s_1=s_2=2$, ${r'}=2$, $A=1$, ${\mu_j}=1$, and ${p'}_j=1/3$.
Each original cell has shape one and probability $1/3$;
their sum has shape two and the same probability.
Its variance is $1+1/2=3/2$, equal to the sum of the original variances.
These two examples separately check the approximation and the exact case.

\subsection{Complexity of assembly, solves, and requested output}
Let $d=GP+R$. The following are dense worst-case algebraic orders,
not timings, speedups, or claims about a particular implementation.
Spline sparsity, caching, and the number of requested contrasts can
substantially change practical costs.
\begin{center}\small
\begin{tabularx}{\textwidth}{T l T}
\toprule
Operation & Dense order & Qualification\\
\midrule
Full $d\times d$ factorization & $O(d^3)$ &
Storage $O(d^2)$, before any map output\\
Factor separate $D_g$ & $O(GP^3)$ &
Storage $O(GP^2)$ for group factors\\
Solve $D_g\mathcal T_g=F_g$ & $O(GP^2R)$ &
Reuses the group factors\\
Form and factor ${\mathcal S}$ & $O(GPR^2+R^3)$ &
Border storage $O(GPR+R^2)$\\
$k$ selected right-hand sides & $O(k[GP^2+GPR+R^2])$ &
After factorization; no full inverse needed\\
Spatial weighted Gram assembly & $O(GNP^2)$ &
Reusable group terms for separable families\\
Study covariate moment assembly & $O(MR^2)$ &
First/second moments for Poisson and matched NB\\
All group map diagonals & $O(GNP^2+GNPR+GNR^2)$ &
Dense triangular solves; chunk voxels to limit memory\\
\bottomrule
\end{tabularx}
\end{center}
Clustered projection also uses study-specific scalar weights and
rank-one group corrections, as in \eqref{eq:cluster-blocks}.
Aggregate NB global blocks use the inner $R$-dimensional moments and
outer voxel sums, not an original independent-cell Hessian.
The displayed assembly orders do not include fitting iterations,
special-function costs, preprocessing, bootstrap replication, or
the cost of constructing a roughness matrix.
An efficient factorization changes neither the model nor the validity
conditions of its covariance approximation.

\subsection{Audit checklist and limits of this appendix}
For manual or implementation review, check in this order:
\begin{enumerate}
 \item Fix the paper's $X,Z,\beta_g,\gamma$ dimensions and remove
 unidentified directions. State the observation model and data margins.
 \item Fix original versus effective dispersion during each derivative.
 Keep factorial and gamma terms whenever their derivatives or absolute
 likelihood interpretation require them.
 \item Verify scalar signs and every nonlinear inner-map term.
 For aggregate NB compare original and log-coordinate formulas;
 for clustered NB compare log-mean projection and total-mean differentiation.
 \item Assemble $D_g,F_g,E$ and compare selected full-system and
 Schur solves, including the signed symmetry and residual checks.
 \item Declare observed or expected information, fixed or joint nuisance
 inference, and frequentist or posterior penalty interpretation.
 \item Form estimable contrasts with the full required covariance and
 qualify the null, asymptotic reference, resampling law, and multiplicity rule.
\end{enumerate}
All derivations are contained in this single source.
The algebra was adapted from the repository's
\texttt{cbmr\_expanded\_appendix.tex} and
\texttt{expanded\_covariance.tex} solely as provenance, not as required reading
or generated inputs. No claim is made that every line has been verified
in a proof assistant, that every optional extension is implemented,
or that the unavailable original supplement has been reconstructed.
Exact differentiation of a specified likelihood, exact block algebra,
working-model accuracy, numerical reliability, and asymptotic coverage
are distinct claims and must be assessed separately.

\subsection{Editable source and independent build}
The editable source is \texttt{docs/imag\_a\_1057\_appendix.tex}.
From the repository root:
\begin{verbatim}
make imag-appendix
\end{verbatim}
The command \texttt{make paper-appendices} builds this document and
the separate lesion-manuscript appendix. Neither command replaces
the supplied papers or the combined expanded appendix.

\begin{thebibliography}{9}
\bibitem{yu2025}
Y. Yu, L. D. Hill-Bowen, M. C. Riedel, et al.
\newblock CBMR: Coordinate-based meta-regression for group and covariate inference.
\newblock \emph{Imaging Neuroscience}, 3, 2025.
\newblock DOI: \href{https://doi.org/10.1162/IMAG.a.1057}{10.1162/IMAG.a.1057}.
\newblock Supplied reference: \texttt{docs/imag.a.1057.pdf}.
\end{thebibliography}
\end{document}
